\chapter{Competition creates research}\label{ch:research}

\begin{chapterplan}
\textbf{The question.} When the informed investor must pay for its information, does a stronger incumbent still
bring informed trade and more entry?

After this chapter you can:
\begin{itemize}
\item explain in words why research raises the posterior after a purchase, and why this lowers what each informed
purchase earns;
\item derive the top posterior \(M_\lambda\), the bottom posterior \(m_\lambda\) and the recruitment intensity
\(\lambda_E(\tau)\) at the prior \(1/2\);
\item derive the research condition \(2a\lambda=\DT m_\lambda-k+(\rho\DT m_\lambda-k)^+\) and show that it has one
root \(\lambda_R(\DT)\) when \(\DT>2k\);
\item state Proposition~4 with its tie rules and prove each part, with the entry formula
\(\rho+(1-\rho)[\beta+(1/2-\beta)\lambda]\);
\item read Corollary~3 and its figure: where zero research, recruiting research and floor-entry research are
equilibria in the example.
\end{itemize}
\end{chapterplan}

% =====================================================================
\section{Research in words}\label{sec:research-words}

In the earlier chapters, the informed investor filled the order slot with a fixed probability \(\lambda\). The paper
calls \(\lambda\) the presence probability. This chapter lets the investor choose it. The investor is a sector
specialist. Before the order opportunity, it can study the target's customers, technology and demand. This work may
reveal the challenger's class. If it does, the investor fills the slot with an order. If it does not, a liquidity
trader fills the slot.

So research is the presence probability of the earlier chapters, chosen at a cost. Each success puts an informed
order in the slot. More research has two effects on what the market sees.
\begin{itemize}
\item \emph{More informed purchases.} After a success in class \(H\), the investor buys. So a purchase becomes more
frequent in class \(H\).
\item \emph{Fewer liquidity orders.} The slot has room for one order. A success takes the place of a liquidity
order, in both classes.
\end{itemize}
Both effects make a purchase better news for class \(H\), so the posterior after a purchase rises. Recall from
Chapter~\ref{ch:trading} that an informed purchase earns entry times the spread times the market's remaining weight
on class \(L\). More research shrinks that remaining weight. Each order hides less, so each success earns less. The
paper says that research reduces the camouflage of the investor's own orders.

A wider target-payoff spread works in the other direction. It raises what an informed order earns at every
posterior. So a wider spread makes research pay more, and the investor researches more. Chapter~\ref{ch:trading}
showed that a stronger incumbent never narrows the spread. This is the route from competition to research.

Research also matters for entry. A high-cost challenger prepares only when its belief reaches the preparation
threshold \(\tau\). With more research, the posterior after a purchase climbs higher. Above some research level, a
purchase recruits a high-cost challenger. The chapter finds when the investor chooses such a level, and what then
happens to entry and to ownership.

One more feature shapes every result. Market makers and the challenger do not see how much research the investor
does. They price and decide at the research level that they expect. In equilibrium, this expectation is correct.

\begin{analogy}[a tipster who chooses how hard to look]
This picture is the guide's own. A betting exchange takes one bet on a horse each day. A tipster can study the
horse's fitness before the bet. If the study works, the tipster places the day's bet, in the right direction. If
it fails, a casual punter places a random bet. The exchange sets fair odds after the bet, and it does not see who
placed it.

If the tipster studies harder, more bets come from the tipster and fewer from casual punters. The exchange expects
this, so it moves the odds more after each bet. Each informed bet then earns less. A larger gap between the payouts
of a fit and an unfit horse makes each informed bet earn more, so the tipster studies harder.

Where the analogy breaks:
\begin{itemize}
\item In the model, the price changes what a share is worth, because it decides whether a high-cost challenger
prepares. Odds do not decide whether a race takes place.
\item The investor earns on its class news only through entry, and it pays a cost \(k\) for each order.
\item Research has a cap \(\lb\) below one: some facts about the match appear only through the challenger's own
diligence.
\end{itemize}
\end{analogy}

\begin{keyidea}
Research is the presence probability chosen at a cost. The value of a success falls with research, because research
removes camouflage, and it rises with the spread. So a wider spread brings more research, and enough research lets a
purchase recruit a high-cost challenger.
\end{keyidea}

% =====================================================================
\section{The set-up}\label{sec:research-setup}

This section recalls the symbols of the chapter and adds research.
\begin{itemize}
\item The challenger's class is \(j\in\{H,L\}\). In this chapter the prior is fixed: \(\Pr(j=H)=1/2\).
Proposition~4 and every result on research use this prior.
\item The preparation cost is \(C\in\{0,c_H\}\), zero with probability \(\rho\in(0,1)\), and \(c_H>0\). A zero-cost
challenger prepares after every price. Entry after the order \(x\) is \(e(x)=\rho+(1-\rho)\,\zeta(P(x))\), where
\(\zeta\) is the probability that a high-cost challenger prepares. So \(e(x)\ge\rho\), the entry floor.
\item A liquidity trader sells, abstains or buys with probabilities \(\beta\), \(1-2\beta\) and \(\beta\). The
research results take \(0<\beta<1/2\) (Appendix A.5 of the paper). A nonzero order costs the investor \(k>0\).
\item \(\DT=t_H-t_L\) is the target-payoff spread. \(g_j\) is the gross acquisition profit of class \(j\),
\(\gamma=g_H-g_L\), and \(\tau=(c_H-g_L)/\gamma\) is the preparation threshold \eqref{eq:beliefs-threshold}. The
research results take ranked classes, a continuous incumbent law, \(\DT>0\), \(\gamma>0\) and \(\tau>1/2\).
\end{itemize}

\paragraph{Research.} Before it learns the class, the investor chooses a research intensity
\(\lambda\in[0,\lb]\), with \(0<\lb<1\), and pays \((a/2)\lambda^2\), with \(a>0\). Research succeeds with probability
\(\lambda\), independently of the class, the costs and the incumbent's value. A success reveals the class, and the
investor then fills the slot with an order. After a failure, a liquidity trader fills the slot.

So the slot works as in Chapter~\ref{ch:model}, with presence probability \(\lambda\). A liquidity order in one
direction has probability \(b=(1-\lambda)\beta\), a liquidity abstention has probability
\(b_0=(1-\lambda)(1-2\beta)\), and \(u=\lambda/b\). The order \(x\) has probability
\(f_j(x)=(1-\lambda)\eta(x)+\lambda\sigma_j(x)\) in class \(j\), by \eqref{eq:model-orderprob}. Here \(\sigma_j\) is
the investor's order distribution after a success in class \(j\).

\paragraph{Why a failure passes the slot on.} The paper gives the reason in one sentence: after a failure, the
investor knows no more than the market, so an order would lose at least its cost \(k\). The guide fills in the step.
After a failure, the investor's belief is the prior \(1/2\). By \eqref{eq:beliefs-terminal}, a share pays
\(t_0+e(q)w_j\) in expectation in class \(j\), where \(w_j=t_j-t_0\) and \(w(\mu)=t_L-t_0+\DT\mu\).

\begin{derivation}[an order without information]
For a purchase:
\begin{align*}
&\tfrac12\bigl[t_0+e(1)w_H\bigr]+\tfrac12\bigl[t_0+e(1)w_L\bigr]-P(1)-k\\
&\quad=t_0+e(1)\,w(1/2)-\bigl[t_0+e(1)\,w(\mu(1))\bigr]-k \why{definition of \(w\); price formula \eqref{eq:beliefs-price}}\\
&\quad=e(1)\,\DT\,\bigl[\tfrac12-\mu(1)\bigr]-k \why{\(w(\mu)=t_L-t_0+\DT\mu\)}
\end{align*}
For a sale, the same steps with the sign reversed give \(e(-1)\,\DT\,[\mu(-1)-\tfrac12]-k\).
\end{derivation}

In a candidate with no wrong-signed orders, \(\mu(-1)\le1/2\le\mu(1)\) (Lemma A.13(d) below). So both brackets are at
most zero, and each uninformed order earns at most \(-k\). The investor therefore leaves the slot to a liquidity
trader.

\paragraph{Why a cap, and why the marginal cost rises.} The cap \(\lb\) is the probability that the investor's work
can reveal the class at all before the challenger decides. Some facts about the match become visible only through
the challenger's own diligence. The marginal cost of research is the derivative of the cost,
\[
\frac{d}{d\lambda}\Bigl[\frac a2\lambda^2\Bigr]=a\lambda .
\]
It rises with \(\lambda\). The paper's reason: the specialist uses cheap sources first, such as its sector coverage,
and costly ones later, such as expert calls. The quadratic cost is the simplest form whose marginal cost rises.

\paragraph{What the market sees.} Market makers and the challenger do not observe the research choice. They use a
conjectured intensity \(\hat\lambda\). Bayes' rule, prices and the entry rule all use \(\hat\lambda\). The paper
defines an equilibrium with research as follows (Appendix A.5 of the paper).
\begin{itemize}
\item Conditions (E1) to (E6) of Chapter~\ref{ch:model} hold at the intensity \(\hat\lambda\).
\item Condition (E5) holds also at \(\hat\lambda=0\). A deviation in research makes a success possible, and the order
after it must be optimal.
\item (E7), \emph{research optimality}. Hold the prices \(P\) and the preparation rule \(\zeta\) fixed. Then
\(\hat\lambda\) is the research choice of a best response among all joint deviations in research and orders.
\end{itemize}
So in equilibrium the conjectured intensity is the investor's best research choice at the equilibrium prices and
entry rule.

\begin{figure}[ht]
\centering
\begin{tikzpicture}[x=1cm,y=1cm,every node/.style={font=\small}]
\draw[rulegrey,thick,-{Stealth[length=2mm]}] (0,0) -- (15.4,0);
\foreach \x/\t in {
  0.9/{The investor chooses research \(\lambda\). No one else sees it.},
  4.0/{Success with probability \(\lambda\): the investor learns the class.},
  7.1/{The slot: the investor's order after a success, a liquidity order after a failure.},
  10.2/{Market makers see the order and set the price.},
  13.3/{The challenger sees the price and its cost, then prepares or not. The auction follows.}}
{
  \draw[ink] (\x,0.15) -- (\x,-0.15);
  \node[anchor=north,text width=2.8cm,align=center] at (\x,-0.25) {\t};
}
\end{tikzpicture}
\caption{The timing with research. Time runs from left to right. Research ends before the order opportunity, and the
class is fixed before research starts.}
\label{fig:research-timeline}
\end{figure}

\begin{watchout}
Keep three features of the set-up.
\begin{itemize}
\item \emph{The prior is \(1/2\).} Proposition~4, Corollary~3 and the lemmas of this chapter fix it. Do not carry their
formulas to another prior.
\item \emph{The domain of \(\beta\).} The research results take \(0<\beta<1/2\). The bound \(\beta\le1/3\) enters only
the bound on the zero-order posterior in Lemma A.13(a) and the slots of Proposition A.1.
\item \emph{Conjecture and choice.} A deviation in research faces fixed prices and entry. A comparison across
incumbents solves the equilibrium again in each economy, research included.
\end{itemize}
\end{watchout}

% =====================================================================
\section{The posteriors at a research level}\label{sec:research-posteriors}

Research raises the top posterior that a purchase can reach. At intensity \(\lambda\), the paper writes \(M_\lambda\)
for the top posterior and \(m_\lambda\) for the bottom posterior. These are the posterior bounds of
Chapter~\ref{ch:beliefs} at the prior \(1/2\), with presence probability \(\lambda\).

\begin{derivation}[the bounds at the prior \(1/2\)]
Start from \eqref{eq:beliefs-bounds}, \(M(\pi)=\pi(1+u)/[\pi(1+u)+1-\pi]\) and \(m(\pi)=\pi/[\pi+(1-\pi)(1+u)]\).
\begin{align*}
M(1/2)&=\frac{\tfrac12(1+u)}{\tfrac12(1+u)+\tfrac12} \why{\(\pi=1-\pi=1/2\)}\\
&=\frac{1+u}{2+u} \why{multiply numerator and denominator by \(2\)}\\[4pt]
m(1/2)&=\frac{\tfrac12}{\tfrac12+\tfrac12(1+u)}=\frac{1}{2+u} \why{the same steps}
\end{align*}
The two add to one: \((1+u)/(2+u)+1/(2+u)=(2+u)/(2+u)=1\). So \(m(1/2)=1-M(1/2)\).
\end{derivation}

Next, put in \(u=\lambda/b\) with \(b=(1-\lambda)\beta\).

\begin{derivation}[the bottom posterior at intensity \(\lambda\)]
\begin{align*}
m_\lambda&=\frac{1}{2+\lambda/b} \why{\(u=\lambda/b\)}\\
&=\frac{b}{2b+\lambda} \why{multiply numerator and denominator by \(b>0\)}\\
&=\frac{(1-\lambda)\beta}{2(1-\lambda)\beta+\lambda} \why{\(b=(1-\lambda)\beta\)}\\
&=\frac{\beta(1-\lambda)}{2\beta+(1-2\beta)\lambda} \why{expand \(2(1-\lambda)\beta\); collect the terms in \(\lambda\)}
\end{align*}
\end{derivation}

This is the paper's formula, with \(M_\lambda=1-m_\lambda\):
\begin{equation}\label{eq:research-bounds}
m_\lambda=\frac{\beta(1-\lambda)}{2\beta+(1-2\beta)\lambda},\qquad
M_\lambda=1-m_\lambda=\frac{b+\lambda}{2b+\lambda}.
\end{equation}
The second form of \(M_\lambda\) follows from \(1-b/(2b+\lambda)=(b+\lambda)/(2b+\lambda)\). It is the posterior after a
purchase when class \(H\) buys after every success and class \(L\) never buys: a purchase has probability \(b+\lambda\)
in class \(H\) and \(b\) in class \(L\) (Lemma A.1(d)).

\paragraph{The top posterior rises with research.} There are two ways to see it.
\begin{itemize}
\item \emph{Through \(u\).} In \(u=\lambda/[(1-\lambda)\beta]\), the numerator rises with \(\lambda\) and the positive
denominator falls. So \(u\) rises strictly. Then \(m_\lambda=1/(2+u)\) falls strictly, and \(M_\lambda=1-m_\lambda\)
rises strictly.
\item \emph{Through the derivative.} The slope formulas of Section~\ref{sec:research-sale} need it.
\end{itemize}

\begin{derivation}[the slope of \(m_\lambda\)]
Write \(m_\lambda=N/D\) with \(N=\beta(1-\lambda)\) and \(D=2\beta+(1-2\beta)\lambda\), so \(N'=-\beta\) and
\(D'=1-2\beta\).
\begin{align*}
m_\lambda'&=\frac{N'D-ND'}{D^2} \why{quotient rule}\\
&=\frac{-\beta\,[2\beta+(1-2\beta)\lambda]-\beta(1-\lambda)(1-2\beta)}{D^2} \why{substitute}\\
&=\frac{-\beta\,[2\beta+(1-2\beta)\lambda+(1-2\beta)-(1-2\beta)\lambda]}{D^2} \why{factor \(-\beta\); expand}\\
&=\frac{-\beta\,[2\beta+1-2\beta]}{D^2}=-\frac{\beta}{[2\beta+(1-2\beta)\lambda]^2} \why{the \(\lambda\) terms cancel}
\end{align*}
\end{derivation}

The slope is negative at every \(\lambda\), so \(m_\lambda\) falls and \(M_\lambda\) rises. Test the extremes. At
\(\lambda=0\), the formula gives \(m_0=\beta/(2\beta)=1/2=M_0\): with no research, every order comes from a liquidity
trader, and a purchase carries no news. As \(\lambda\) moves toward one, \(m_\lambda\) moves toward zero and
\(M_\lambda\) toward one. The cap \(\lb<1\) keeps \(M_\lambda\) below one.

\begin{intuition}[two reasons in one formula]
Write \(M_\lambda=(b+\lambda)/(2b+\lambda)\). The likelihood ratio of a purchase is \((b+\lambda)/b=1+u\). Research
raises the numerator \(b+\lambda\), because it adds informed purchases in class \(H\). It lowers \(b\), because it
removes liquidity purchases in both classes. The two changes push the ratio up together. This is why more research
leaves each purchase with less camouflage.
\end{intuition}

\paragraph{The recruitment intensity.} When \(1/2<\tau<1\), the paper defines the recruitment intensity
\(\lambda_E(\tau)\) as the research level at which a full purchase just reaches the threshold: \(M_\lambda=\tau\).
Since \(M_\lambda\) rises strictly, a full purchase reaches the threshold exactly when \(\lambda\ge\lambda_E(\tau)\).
The algebra is shortest through \(u\).

\begin{derivation}[the recruitment intensity]
First solve for \(u\):
\begin{align*}
\frac{1+u}{2+u}=\tau
&\iff 1+u=2\tau+\tau u \why{multiply by \(2+u>0\)}\\
&\iff u(1-\tau)=2\tau-1 \why{collect the terms in \(u\)}\\
&\iff u=\frac{2\tau-1}{1-\tau} \why{divide by \(1-\tau>0\)}
\end{align*}
Next solve \(u=\lambda/[(1-\lambda)\beta]\) for \(\lambda\):
\begin{align*}
u\beta(1-\lambda)&=\lambda \why{multiply by \((1-\lambda)\beta>0\)}\\
u\beta&=\lambda\,(1+u\beta) \why{add \(u\beta\lambda\) to both sides}\\
\lambda&=\frac{u\beta}{1+u\beta} \why{divide by \(1+u\beta>0\)}
\end{align*}
Put in the value of \(u\):
\begin{align*}
\lambda_E(\tau)&=\frac{\beta(2\tau-1)/(1-\tau)}{1+\beta(2\tau-1)/(1-\tau)} \why{substitute \(u\)}\\
&=\frac{\beta(2\tau-1)}{(1-\tau)+\beta(2\tau-1)} \why{multiply numerator and denominator by \(1-\tau\)}\\
&=\frac{\beta(2\tau-1)}{\beta+(1-2\beta)(1-\tau)} \why{\((1-\tau)+\beta(2\tau-1)=\beta+(1-2\beta)(1-\tau)\)}
\end{align*}
\end{derivation}

To check the last identity, expand its right side: \(\beta+(1-\tau)-2\beta(1-\tau)=(1-\tau)+\beta(2\tau-1)\). The result
is the paper's formula:
\begin{equation}\label{eq:research-lamE}
\lambda_E(\tau)=\frac{\beta(2\tau-1)}{\beta+(1-2\beta)(1-\tau)},\qquad m_{\lambda_E(\tau)}=1-\tau .
\end{equation}

Three facts follow.
\begin{itemize}
\item \emph{Domain.} The map \(u\mapsto u\beta/(1+u\beta)\) sends every positive \(u\) into \((0,1)\), and
\(u=(2\tau-1)/(1-\tau)\) is positive exactly when \(1/2<\tau<1\). So \(0<\lambda_E(\tau)<1\) on this range. If
\(\tau\ge1\), no research level recruits, because \(M_\lambda<1\).
\item \emph{Slope.} By the chain rule, \(d\lambda/du=\beta/(1+u\beta)^2\) and \(du/d\tau=[2(1-\tau)+(2\tau-1)]/(1-\tau)^2
=1/(1-\tau)^2\). Their product is \(\beta/[(1-\tau)(1+u\beta)]^2\), and \((1-\tau)(1+u\beta)=(1-\tau)+\beta(2\tau-1)\). So
\[
\frac{d\lambda_E}{d\tau}=\frac{\beta}{[\beta+(1-2\beta)(1-\tau)]^2}>0 .
\]
A higher threshold needs more research before a purchase recruits.
\item \emph{The cap of the example.} The example sets \(\lb=u\beta/(1+u\beta)\) with \(1+u=\eu\). This is the same
map from \(u\) to \(\lambda\), so research at the cap reproduces the example's slot, with \(M_{\lb}=\eu/(1+\eu)\).
\end{itemize}

\begin{checkyourself}\label{q:research-lamE}
Derive \(\lambda_E(\tau)\) a second way: solve \(m_\lambda=1-\tau\) for \(\lambda\) directly from
\eqref{eq:research-bounds}. Then show that \(\lambda_E\) rises with \(\beta\) at a fixed \(\tau\in(1/2,1)\), and explain
the direction in words.
\end{checkyourself}

% =====================================================================
\section{The value of a success and the research condition}\label{sec:research-value}

\subsection{The research choice at fixed prices}

The investor chooses research before it learns the class, and it trades only after a success. The paper first
solves its problem at fixed prices and entry.

\begin{paperresult}{Lemma A.12 (Joint deviations)}
For a candidate satisfying beliefs, competitive pricing, preparation and truthful bidding, let
\(\Psi=\tfrac12\max_qU_H(q)+\tfrac12\max_qU_L(q)\) be its best trading value after a research success. At its fixed
prices and preparation rule, the best payoff of a joint deviation to any \(\tilde\lambda\in[0,\lb]\) is
\(\tilde\lambda\Psi-(a/2)\tilde\lambda^2\). The unique best research choice is \(\min\{\lb,\Psi/a\}\). Thus research
optimality is exactly \(\hat\lambda=\min\{\lb,\Psi/a\}\), under all four tie rules.
\end{paperresult}

In plain words: \(U_j(q)\) is the payoff of the order \(q\) in class \(j\) at the candidate's prices and entry, as in
Chapter~\ref{ch:trading}. \(\Psi\) is what a success is worth: the class is \(H\) or \(L\) with probability \(1/2\)
each, and the investor then places its best order. The research choice equals this value divided by \(a\), cut at the
cap.

\paragraph{Proof, step by step.}
\begin{enumerate}
\item \emph{Prices do not move.} A joint deviation \((\tilde\lambda,\sigma')\) changes the true order probabilities to
\((1-\tilde\lambda)\eta+\tilde\lambda\sigma'_j\). It does not change the prices or the preparation rule, because
market makers and the challenger use the candidate's \(\hat\lambda\) and \(\sigma_j\).
\item \emph{The expected payoff.} With probability \(1-\tilde\lambda\), research fails, the investor does not trade,
and it earns zero. With probability \(\tilde\lambda\), it learns the class and places an order. So
\[
\Pi(\tilde\lambda,\sigma')=\tilde\lambda\Bigl[\tfrac12\sum_q\sigma'_H(q)U_H(q)+\tfrac12\sum_q\sigma'_L(q)U_L(q)\Bigr]
-\frac a2\tilde\lambda^2 .
\]
\item \emph{The best orders.} Each sum is an average of the payoffs \(U_j(q)\), so it is at most \(\max_qU_j(q)\). The
bracket is at most \(\Psi\), with equality exactly when each \(\sigma'_j\) puts weight only on best orders. So the best
payoff at \(\tilde\lambda\) is \(\tilde\lambda\Psi-(a/2)\tilde\lambda^2\).
\item \emph{The best research.} The derivative in \(\tilde\lambda\) is \(\Psi-a\tilde\lambda\), and the second
derivative is \(-a<0\), so the function is strictly concave. Also \(\Psi\ge0\), because \(U_j(0)=0\).
\begin{itemize}
\item If \(\Psi\ge a\lb\), the derivative is positive on \([0,\lb)\), and the maximizer is \(\lb\).
\item If \(\Psi<a\lb\), the derivative is zero at \(\Psi/a\in[0,\lb)\), and this is the maximizer.
\end{itemize}
Together, the maximizer is \(\min\{\lb,\Psi/a\}\).
\item \emph{No tie.} Strict concavity makes the maximizer unique. It also excludes a mixture over research levels: a
mixture puts weight on levels with a strictly lower payoff.
\end{enumerate}
The tie rules concern indifferent orders and indifferent preparation. Neither enters this problem, so the lemma holds
under all four.

\paragraph{Why the investor takes the posterior as given.} In \(\Psi\), the payoffs \(U_j(q)\) contain the posteriors
\(\mu(q)\) and entry \(e(q)\) of the candidate. Market makers compute these with the conjectured \(\hat\lambda\). So,
to the investor, \(\Psi\) is a fixed number, and its problem is a simple concave quadratic. Equilibrium then asks that
the choice equal the conjecture. The research condition below therefore evaluates the posteriors at the chosen
intensity.

\begin{intuition}[the camouflage that the investor does not count]
This reading is the guide's own. Because no one sees its research, the investor cannot change the prices by its
choice. So, at the margin, it does not count the camouflage that one more unit of research removes. The market makers
count it, through their conjecture: a higher conjectured intensity makes them read each purchase as better
informed. The loss of camouflage therefore appears in the research condition only through the conjecture.
\end{intuition}

\subsection{What a success is worth}

Take the equilibrium of Proposition~4: class \(H\) buys after a success, and a purchase recruits a high-cost
challenger. Its value of a success has one term for each class.

\paragraph{Class \(H\).} Class \(H\) buys for sure, so \(\mu(1)=M_\lambda\). The purchase recruits, so \(e(1)=1\). By
\eqref{eq:trading-returns},
\begin{align*}
U_H(1)&=e(1)\,\DT\,[1-\mu(1)]-k \why{\eqref{eq:trading-returns}}\\
&=\DT\,(1-M_\lambda)-k \why{\(e(1)=1\), \(\mu(1)=M_\lambda\)}\\
&=\DT\,m_\lambda-k \why{\(1-M_\lambda=m_\lambda\)}
\end{align*}
Abstention earns zero, and a sale earns at most \(-k\) (Lemma A.3(a)). So the best payoff of class \(H\) is
\(\DT m_\lambda-k\) when this is positive.

\paragraph{Class \(L\).} A sale never recruits: its posterior is at most \(1/2<\tau\). So class \(L\) faces the problem
of Chapter~\ref{ch:trading} at entry \(\rho\). The paper solves it in a lemma.

\begin{paperresult}{Lemma A.15 (The low class's sale)}
Let \(0<\lambda\le\lb\), let class \(H\) buy with probability one after a success, and suppose an indifferent investor
may mix over its best orders. By Lemma A.8 at presence \(\lambda\) and prior \(1/2\), the unique low-class sale rate
consistent with its best orders is
\[
s_-=\lambda\sigma_L(-1)=\min\{\lambda,\ b(\rho\DT/k-2)^+\}.
\]
The remaining low-class probability is on abstention, and
\[
\mu(-1)=\frac{b}{2b+s_-},\qquad \mu(0)=\frac{b_0}{2b_0+\lambda-s_-},\qquad \mu(1)=M_\lambda.
\]
Both nonpurchase orders have entry \(\rho\), under either challenger tie rule. Class \(L\) abstains when
\(\rho\DT\le2k\); it mixes when \(\rho\DT m_\lambda<k<\rho\DT/2\), with \(\mu(-1)=k/(\rho\DT)\); and it sells with
probability one when \(\rho\DT m_\lambda\ge k\). At \(\rho\DT=2k\) its sale rate is zero; at \(\rho\DT m_\lambda=k\)
its sale probability is one. Its best net payoff is \(\max_qU_L(q)=(\rho\DT m_\lambda-k)^+\).
\end{paperresult}

Lemma A.8 is the low-class lemma of Chapter~\ref{ch:coexistence} at a fixed presence probability. The proof below
redoes its steps at the prior \(1/2\).

\begin{derivation}[proof of Lemma A.15]
\emph{Step 1: order probabilities.} Class \(L\) never buys (Lemma A.3(a)). Class \(H\) buys after every success. On
the orders \((-1,0,1)\),
\[
f_H=(b,\;b_0,\;b+\lambda),\qquad f_L=(b+s_-,\;b_0+\lambda-s_-,\;b).
\]
\emph{Step 2: posteriors.} At the prior \(1/2\), Bayes' rule reads \(\mu(x)=f_H(x)/[f_H(x)+f_L(x)]\), because the
factor \(1/2\) cancels. This gives the three displayed posteriors.

\emph{Step 3: entry.} \(\mu(-1)=b/(2b+s_-)\le1/2\), and \(\mu(0)\le1/2\) because \(\lambda-s_-\ge0\). Both are below
\(\tau\), so entry after a sale and after no order is \(\rho\) under either challenger rule (Lemma A.2(d)).

\emph{Step 4: the sale payoff.}
\begin{align*}
U_L(-1)&=e(-1)\,\DT\,\mu(-1)-k \why{\eqref{eq:trading-returns}}\\
&=\frac{\rho\DT\,b}{2b+s_-}-k \why{\(e(-1)=\rho\); Step 2}
\end{align*}
The denominator rises with \(s_-\), so \(U_L(-1)\) falls strictly in \(s_-\). Its two end values are:
\begin{itemize}
\item at \(s_-=0\): \(\rho\DT/2-k\), which is at most zero exactly when \(\rho\DT\le2k\);
\item at \(s_-=\lambda\): \(\rho\DT\,b/(2b+\lambda)-k=\rho\DT m_\lambda-k\), by \eqref{eq:research-bounds}.
\end{itemize}

\emph{Step 5: the three cases.} Under (E5), class \(L\) sells only if the sale is a best order.
\begin{itemize}
\item If \(\rho\DT\le2k\), every positive rate makes the sale lose strictly. So \(s_-=0\), and class \(L\) abstains.
\item If \(\rho\DT m_\lambda\ge k\), every rate below \(\lambda\) makes the sale strictly profitable. So
\(s_-=\lambda\), and class \(L\) sells for sure.
\item Otherwise \(\rho\DT m_\lambda<k<\rho\DT/2\). The sale payoff is positive at \(s_-=0\) and negative at
\(s_-=\lambda\). By the intermediate value theorem and strict monotonicity, it is zero at exactly one rate in
between. An indifferent investor may mix, so class \(L\) sells at this rate.
\end{itemize}
\emph{Step 6: the mixing rate.} Set the sale payoff to zero:
\begin{align*}
\frac{\rho\DT\,b}{2b+s_-}=k
&\iff 2b+s_-=\frac{\rho\DT\,b}{k} \why{multiply by \((2b+s_-)/k\)}\\
&\iff s_-=b\Bigl(\frac{\rho\DT}{k}-2\Bigr) \why{subtract \(2b\)}
\end{align*}
At this rate, \(\mu(-1)=k/(\rho\DT)\), from the first line. The formula \(\min\{\lambda,b(\rho\DT/k-2)^+\}\) collects the
three cases: the positive part is zero when \(\rho\DT\le2k\), and \(b(\rho\DT/k-2)\ge\lambda\) is the same as
\(\rho\DT/k\ge2+u\), that is \(\rho\DT m_\lambda\ge k\).

\emph{Step 7: the best payoff.} When class \(L\) abstains or mixes, its best payoff is zero. When it sells for sure, it
is \(\rho\DT m_\lambda-k\ge0\). In the first two cases \(\rho\DT m_\lambda<k\), because \(m_\lambda<1/2\). So the best
payoff is \((\rho\DT m_\lambda-k)^+\) in all three cases.
\end{derivation}

\subsection{The research condition}

Add the two classes with weight \(1/2\) each:
\[
\Psi=\tfrac12\bigl[\DT m_\lambda-k\bigr]+\tfrac12\bigl(\rho\DT m_\lambda-k\bigr)^+ .
\]
Lemma A.12 gives research \(\lambda=\min\{\lb,\Psi/a\}\). At an interior intensity, \(\lambda<\lb\):

\begin{derivation}[the research condition]
\begin{align*}
\lambda&=\frac{\Psi}{a} \why{Lemma A.12, below the cap}\\
a\lambda&=\tfrac12\bigl[\DT m_\lambda-k+(\rho\DT m_\lambda-k)^+\bigr] \why{multiply by \(a\); substitute \(\Psi\)}\\
2a\lambda&=\DT m_\lambda-k+(\rho\DT m_\lambda-k)^+ \why{multiply by \(2\)}
\end{align*}
\end{derivation}

This is equation (10) of the paper:
\begin{equation}\label{eq:research-cond}
2a\lambda=\DT\,m_\lambda-k+\bigl(\rho\DT\,m_\lambda-k\bigr)^+ .
\end{equation}
The left side is twice the marginal cost of research. The right side is twice the value of a success. At the cap,
Lemma A.12 asks instead that the left side be at most the right side at \(\lambda=\lb\).

\begin{intuition}[why a factor of two]
A success reveals class \(H\) only half the time. Then a recruiting purchase earns \(\DT m_\lambda-k\). The other
half of the time, class \(L\) earns its best payoff. Each success is worth the average of the two, and the investor
pushes research until the marginal cost \(a\lambda\) equals this average.
\end{intuition}

\paragraph{The two sides.} Write the right side minus the left side as the residual
\[
\mathcal R(\lambda)=\DT m_\lambda-k+(\rho\DT m_\lambda-k)^+-2a\lambda .
\]
\begin{itemize}
\item The left side \(2a\lambda\) rises strictly.
\item The right side falls strictly. The term \(\DT m_\lambda\) falls strictly, because \(m_\lambda\) does and
\(\DT>0\). The term \(\rho\DT m_\lambda-k\) also falls, and the positive part never falls when its argument rises
(Section~\ref{sec:toolkit-monotone}). So the positive part never rises with \(\lambda\).
\end{itemize}
So \(\mathcal R\) is continuous and strictly decreasing. In the paper's words, the right side falls because more
research removes camouflage. The fall is a fall in the conjectured intensity: market makers who expect more research
read each purchase as better informed.

\paragraph{One root when \(\DT>2k\).} Evaluate \(\mathcal R\) at the two ends of \([0,1]\). The formula
\eqref{eq:research-bounds} extends to \(\lambda=1\), with \(m_1=0\). The point \(\lambda=1\) lies outside the model; it
serves only the intermediate value theorem.
\begin{align*}
\mathcal R(0)&=\tfrac12\DT-k+(\tfrac12\rho\DT-k)^+\ge\tfrac12\DT-k \why{\(m_0=1/2\); a positive part is \(\ge0\)}\\
\mathcal R(1)&=-k+(-k)^+-2a=-k-2a<0 \why{\(m_1=0\); \((-k)^+=0\)}
\end{align*}
If \(\DT>2k\), then \(\mathcal R(0)>0>\mathcal R(1)\). The intermediate value theorem and strict decrease give
exactly one root in \((0,1)\). The paper calls it \(\lambda_R(\DT)\).

If \(\DT\le2k\), every positive \(\lambda\) has \(\DT m_\lambda<\DT/2\le k\). The purchase then loses money, and
class \(H\) would not buy. So this equilibrium needs \(\DT>2k\).

\paragraph{Research is \(\min\{\lambda_R,\lb\}\).} Condition (E7) reads \(\hat\lambda=\min\{\lb,\Psi(\hat\lambda)/a\}\),
where \(\Psi(\hat\lambda)\) is the value of a success at the conjectured intensity. The guide shows that this is the
same as \(\hat\lambda=\min\{\lambda_R,\lb\}\). Note that \(\Psi(\lambda)/a\ge\lambda\) exactly when
\(\mathcal R(\lambda)\ge0\), with equality in both at once.
\begin{itemize}
\item \emph{Suppose \(\hat\lambda<\lb\).} Then the minimum is \(\Psi/a\), so \(\hat\lambda=\Psi(\hat\lambda)/a\), and
\(\mathcal R(\hat\lambda)=0\). So \(\hat\lambda=\lambda_R<\lb\), and \(\hat\lambda=\min\{\lambda_R,\lb\}\).
\item \emph{Suppose \(\hat\lambda=\lb\).} Then \(\Psi(\lb)/a\ge\lb\), so \(\mathcal R(\lb)\ge0\). Since \(\mathcal R\)
is strictly decreasing with root \(\lambda_R\), this gives \(\lb\le\lambda_R\). Again \(\hat\lambda=\min\{\lambda_R,\lb\}\).
\item \emph{Conversely.} If \(\lambda_R<\lb\), then \(\mathcal R(\lambda_R)=0\), so \(\Psi/a=\lambda_R<\lb\), and (E7)
holds at \(\lambda_R\). If \(\lambda_R\ge\lb\), then \(\mathcal R(\lb)\ge0\), so \(\Psi(\lb)/a\ge\lb\), and (E7) holds
at \(\lb\).
\end{itemize}
Research never has a tie, by the last step of Lemma A.12.

\paragraph{An explicit root.} While the sale term is zero at the root, \eqref{eq:research-cond} reads
\(2a\lambda=\DT m_\lambda-k\). Multiply by \(D=2\beta+(1-2\beta)\lambda\). The paper does this for the example
(equation (A.16)).
\begin{align*}
2D\Bigl[a\lambda-\tfrac12(\DT m_\lambda-k)\Bigr]
&=2a\lambda D-\DT\beta(1-\lambda)+kD \why{\(D\,m_\lambda=\beta(1-\lambda)\)}\\
&=2a(1-2\beta)\lambda^2+\bigl[4a\beta+k(1-2\beta)+\beta\DT\bigr]\lambda+\beta(2k-\DT) \why{expand; collect powers of \(\lambda\)}
\end{align*}
Call this quadratic \(\Lambda(\lambda)\), with slope \(B=4a\beta+k(1-2\beta)+\beta\DT\) at zero. Its constant
\(\beta(2k-\DT)\) is negative when \(\DT>2k\), and its other coefficients are positive, so it has one positive root.
The quadratic formula, with the numerator and denominator multiplied by \(B+\sqrt{\cdot}\), gives
\[
\lambda_R(\DT)=\frac{2\beta(\DT-2k)}{B+\sqrt{B^2+8a\beta(1-2\beta)(\DT-2k)}} .
\]
The numerator \(\DT-2k\) shows again that research needs \(\DT>2k\).

\begin{figure}[ht]
\centering
\begin{tikzpicture}[x=13cm,y=8.5cm,every node/.style={font=\small}]
\draw[-{Stealth[length=2mm]}] (0,0) -- (0.86,0) node[right] {\(\lambda\)};
\draw[-{Stealth[length=2mm]}] (0,-0.06) -- (0,0.6);
\node[below left] at (0,0) {\(0\)};
\draw[accent,thick,domain=0:0.8,samples=60] plot (\x,{0.3*(1-\x)/(0.6+0.4*\x)-0.15});
\draw[ink,thick,domain=0:0.8,samples=60] plot (\x,{1.3*0.3*(1-\x)/(0.6+0.4*\x)-0.15});
\draw[notegrey,thick,dashed] (0,0) -- (0.8,0.48) node[right,text=ink] {\(2a\lambda\)};
\node[right,text=accent] at (0.8,{0.3*0.2/(0.6+0.32)-0.15}) {\(\DT m_\lambda-k\)};
\node[right] at (0.8,{1.3*0.3*0.2/(0.6+0.32)-0.15}) {\(\DT' m_\lambda-k\)};
\draw[dotted] (0.2857,0) -- (0.2857,0.42);
\node[above] at (0.2857,0.42) {\(\lambda_E(\tau)\)};
\draw[dotted] (0.27,0.162) -- (0.27,0) node[below] {\(\lambda_R\)};
\draw[dotted] (0.338,0.203) -- (0.338,0) node[below] {\(\lambda_R'\)};
\fill[accent] (0.27,0.162) circle (0.25pt);
\fill[ink] (0.338,0.203) circle (0.25pt);
\draw[dotted] (0.62,0) -- (0.62,0.5);
\node[above] at (0.62,0.5) {\(\lb\)};
\end{tikzpicture}
\caption{The research condition while the sale term is zero. The dashed line is twice the marginal cost. Each falling
curve is twice the value of a success, for the spread \(\DT\) and for a wider spread \(\DT'\). Research is where a
curve meets the line. With the narrower spread, the meeting point lies left of \(\lambda_E(\tau)\), so a purchase
falls short of the threshold. With the wider spread it lies right of \(\lambda_E(\tau)\), so a purchase recruits. The
positions are only a picture.}
\label{fig:research-condition}
\end{figure}

\paragraph{What this says.} Research balances twice the marginal cost against twice the value of a success. A wider
spread lifts the falling curve and moves the meeting point to the right (Figure~\ref{fig:research-condition}). Test
the extremes.
\begin{itemize}
\item \emph{A spread at or below \(2k\).} Even with no research at all, a purchase at the posterior \(1/2\) earns at
most its cost. No recruiting research exists.
\item \emph{Cheap research, \(a\) near zero.} The line \(2a\lambda\) lies flat near zero, so the root moves toward the
point where \(\DT m_\lambda=k\). If that point exceeds the cap, research sits at the cap.
\item \emph{Costly research, \(a\) large.} The line is steep, and the root moves toward zero. A purchase then carries
little news.
\end{itemize}

\subsection{The exact recruitment margin}

The purchase recruits when \(M_{\hat\lambda}\ge\tau\), that is, when research is at least \(\lambda_E(\tau)\).
Figure~\ref{fig:research-condition} shows how to read this off the residual: compare the curve and the line at
\(\lambda_E(\tau)\). The paper defines the \emph{exact recruitment margin} as the right side of
\eqref{eq:research-cond} minus its left side at the recruitment intensity. With \(J=\DT(1-\tau)\) and
\(m_{\lambda_E}=1-\tau\) from \eqref{eq:research-lamE},
\begin{equation}\label{eq:research-margin}
\varphi_a^*=\mathcal R(\lambda_E(\tau))=J-k+(\rho J-k)^+-2a\lambda_E(\tau).
\end{equation}
Recall from Section~\ref{sec:trading-path} that \(J\) is the gross return of a purchase that just reaches \(\tau\) and
recruits. So \(\varphi_a^*\) is positive exactly when, at the research level that just lets a purchase recruit, the
value of a success exceeds the marginal cost of research.

\begin{derivation}[recruitment and the margin, when \(1/2<\tau<1\)]
\begin{align*}
M_{\min\{\lambda_R,\lb\}}\ge\tau
&\iff \min\{\lambda_R,\lb\}\ge\lambda_E(\tau) \why{\(M_\lambda\) rises strictly; \(M_{\lambda_E}=\tau\)}\\
&\iff \lambda_R\ge\lambda_E(\tau)\ \text{and}\ \lb\ge\lambda_E(\tau) \why{a minimum is \(\ge x\) iff both terms are}\\
&\iff \varphi_a^*\ge0\ \text{and}\ \lambda_E(\tau)\le\lb \why{\(\mathcal R\) strictly decreasing, zero at \(\lambda_R\)}
\end{align*}
\end{derivation}

A nonnegative margin also forces \(\DT>2k\), so the root exists. Suppose \(\varphi_a^*\ge0\). Then
\(J-k+(\rho J-k)^+\ge2a\lambda_E>0\). If \(J\le k\), then \(\rho J-k<0\), and the left side equals \(J-k\le0\), a
contradiction. So \(J>k\). Since \(1-\tau<1/2\), \(\DT/2>\DT(1-\tau)=J>k\), and \(\DT>2k\).

\begin{watchout}
Read \eqref{eq:research-cond} as research optimality at an \emph{interior} intensity. At the cap the equality becomes
the weak inequality \(2a\lb\le\DT m_{\lb}-k+(\rho\DT m_{\lb}-k)^+\). Also keep the two roles of \(m_\lambda\). The
investor takes it as given, because prices use the conjectured intensity. Equilibrium evaluates it at the chosen
intensity.
\end{watchout}

% =====================================================================
\section{Proposition 4: competition creates research}\label{sec:research-prop4}

Here is the result as the paper states it in the main text. The box copies every hypothesis and the tie rules.

\begin{paperresult}{Proposition 4 (Competition creates research)}
Let the classes satisfy the class ranking and the prior be \(1/2\), and let an indifferent investor mix over its best
orders and an indifferent high-cost challenger prepare. Take incumbent distributions with a positive spread,
\(g_H>g_L\) and \(\tau>1/2\), and consider the equilibrium in which class \(H\) buys after a success and a purchase
recruits a high-cost challenger. The same recruiting outcome is available when an indifferent high-cost challenger
may prepare with any probability and prepares fully at the purchase.
\begin{enumerate}[label=(\alph*),leftmargin=2em]
\item The equilibrium exists if and only if \(\DT>2k\) and \(\tau\le M_\lambda\) at \(\lambda=\min\{\lambda_R(\DT),\lb\}\),
which is then the investor's research.
\item Research rises strictly with the spread until it reaches the cap and then stays at the cap. Entry,
\(\rho+(1-\rho)[\beta+(1/2-\beta)\lambda]\), and entry in class \(H\) rise with research, and entry in class \(L\) falls.
\item When class \(H\) outbids the incumbent, high-class ownership equals half the entry in class \(H\). Along a path on
which class \(H\) outbids every incumbent, ownership does not fall with strength, and it rises strictly where research
is below the cap and the spread rises strictly.
\item Along a compact path of such distributions with continuous spread and threshold, let the spread against the
weakest incumbent lie below the value at which research reaches its cap, let \(r_{\lb}\) be the first strength at
which the spread reaches that value, and let the threshold stay below one up to \(r_{\lb}\). If the exact recruitment
margin rises strictly up to \(r_{\lb}\), from a negative value against the weakest incumbent to a positive value at
\(r_{\lb}\), then the equilibrium exists exactly on an interval of strengths. On it, research rises from the
recruitment intensity to the cap and then stays at the cap, up to the last strength at which the threshold does not
exceed the top posterior at the cap.
\end{enumerate}
\end{paperresult}

A plain reading:
\begin{itemize}
\item \emph{Hypotheses.} Ranked classes, the prior \(1/2\), a positive spread, a positive profit gap and a threshold
above the prior. The appendix setting adds a continuous incumbent law, \(k,a>0\), \(0<\beta<1/2\) and \(0<\lb<1\).
\item \emph{Tie rules.} The default tie rules: an indifferent investor may mix, and an indifferent high-cost challenger
prepares. Under the other challenger rule, the paper selects full preparation after a purchase that reaches the
threshold, and the same outcome exists.
\item \emph{The equilibrium.} The paper calls it \emph{recruiting research}: positive research, and full preparation
after a purchase.
\item \emph{(a)} It exists exactly when the spread exceeds twice the trading cost and a purchase at the chosen research
reaches the threshold. \emph{(b)} More spread brings more research, up to the cap, and more entry, tilted toward
class \(H\). \emph{(c)} Ownership by class \(H\) follows its entry when class \(H\) always wins. \emph{(d)} Along a
path, the equilibrium holds on one interval of strengths.
\end{itemize}

\subsection{Two lemmas behind part (a)}

The proof of part (a) uses two lemmas. The first holds for every order profile.

\begin{paperresult}{Lemma A.13 (Bounds under every profile)}
Let \(0\le\hat\lambda\le\lb<1\). Under all four tie rules:
\begin{enumerate}[label=(\alph*),leftmargin=2em]
\item For every order profile, the nonzero-order posteriors lie in \([m_{\hat\lambda},M_{\hat\lambda}]\), with
\(M_{\hat\lambda}\le M_{\lb}\). If also \(\beta\le1/3\), the same bound holds for the zero-order posterior under every
profile.
\item For every candidate satisfying beliefs, competitive pricing, preparation and truthful bidding, \(P(q)=P(q')\)
implies \(\mu(q)=\mu(q')\), and \(\bar\mu(P(q))=\mu(q)\).
\item Wrong-signed orders satisfy \(U_H(-1)\le-k\) and \(U_L(1)\le-k\); every equilibrium has
\(\sigma_H(-1)=\sigma_L(1)=0\).
\item With no wrong-signed orders, \(\mu(-1)\le1/2\le\mu(1)\). If also \(\sigma_H(1)=1\), then
\(\mu(1)=M_{\hat\lambda}\) and \(\mu(0)\le1/2\).
\end{enumerate}
\end{paperresult}

Write \(\hat b=(1-\hat\lambda)\beta\) and \(\hat b_0=(1-\hat\lambda)(1-2\beta)\). Each part repeats a step of an earlier
chapter at the conjectured intensity.
\begin{enumerate}
\item \emph{Part (a).} A nonzero order has probability between \(\hat b\) and \(\hat b+\hat\lambda\) in each class
(Lemma A.1(b)). At the prior \(1/2\), the posterior then lies between \(\hat b/(2\hat b+\hat\lambda)=m_{\hat\lambda}\)
and \(M_{\hat\lambda}\). The zero order uses \(\hat b_0\) in place of \(\hat b\); its bounds lie inside when
\(\hat b_0\ge\hat b\), that is \(\beta\le1/3\) (Section~\ref{sec:beliefs-zero}). The top bound rises with research, so
\(M_{\hat\lambda}\le M_{\lb}\).
\item \emph{Part (b).} This is Lemma A.2(c): equal prices force equal entry, and with \(e\ge\rho>0\) and \(\DT>0\),
equal posteriors (Section~\ref{sec:beliefs-reveal}).
\item \emph{Part (c).} This is Lemma A.3(a): each wrong-signed payoff is at most \(-k<0=U_j(0)\), so neither investor
rule puts weight on it.
\item \emph{Part (d).} Without wrong-signed orders, \(f_H(1)=\hat b+\hat\lambda\sigma_H(1)\ge\hat b=f_L(1)\) and
\(f_H(-1)=\hat b\le\hat b+\hat\lambda\sigma_L(-1)=f_L(-1)\). A likelihood ratio at least one gives a posterior at least
the prior \(1/2\), and the reverse. With \(\sigma_H(1)=1\), \(f_H(0)=\hat b_0\le\hat b_0+\hat\lambda\sigma_L(0)=f_L(0)\),
so \(\mu(0)\le1/2\), and \(\mu(1)=(\hat b+\hat\lambda)/(2\hat b+\hat\lambda)=M_{\hat\lambda}\).
\end{enumerate}

The second lemma says that positive research needs a purchase that pays strictly.

\begin{paperresult}{Lemma A.14 (Positive research needs a strict full purchase)}
Under either investor tie rule and either challenger tie rule, every equilibrium with \(\hat\lambda>0\) has
\(U_H(1)>0\) and \(\sigma_H(1)=1\). Its posteriors satisfy \(\mu(1)=M_{\hat\lambda}\), \(\mu(0)\le1/2\) and
\(\mu(-1)\le1/2\); entry after a sale or no order is \(\rho\). This holds for every \(0<\beta<1/2\).
\end{paperresult}

\begin{derivation}[proof of Lemma A.14]
\emph{Step 1: some correct order pays.} By Lemma A.12, \(\hat\lambda=\min\{\lb,\Psi/a\}>0\), so \(\Psi>0\). Wrong-signed
orders earn at most \(-k\) and abstention earns zero, so \(U_H(1)>0\) or \(U_L(-1)>0\).

\emph{Step 2: a paying sale implies a paying purchase.} Suppose \(U_L(-1)>0\). Then the sale is the unique best order of
class \(L\), and \(\sigma_L(-1)=1\). With \(\sigma_H(-1)=0\), \(f_H(-1)=\hat b\) and \(f_L(-1)=\hat b+\hat\lambda\), so
\(\mu(-1)=m_{\hat\lambda}\le1/2<\tau\). Hence \(e(-1)=\rho\) and \(U_L(-1)=\rho\DT m_{\hat\lambda}-k>0\). For the
purchase,
\begin{align*}
U_H(1)&=e(1)\,\DT\,[1-\mu(1)]-k \why{\eqref{eq:trading-returns}}\\
&\ge\rho\,\DT\,(1-M_{\hat\lambda})-k \why{\(e(1)\ge\rho\); \(\mu(1)\le M_{\hat\lambda}\), Lemma A.13(a)}\\
&=\rho\DT m_{\hat\lambda}-k>0 \why{\(1-M=m\); the line above}
\end{align*}
So \(U_H(1)>0\) in both cases.

\emph{Step 3: the consequences.} The purchase is the unique best order of class \(H\), so \(\sigma_H(1)=1\) under both
investor rules. Lemma A.13(d) gives the three posteriors. Since \(\tau>1/2\), Lemma A.2(d) gives entry \(\rho\) after a
sale and after no order, under both challenger rules.
\end{derivation}

\subsection{Part (a): when recruiting research exists}

The appendix states part (a) in full.

\begin{paperresult}{Proposition 4(a), full statement}
When an indifferent high-cost challenger may prepare with any probability, select full preparation after a purchase
that reaches the threshold. Recruiting research, an equilibrium with positive research and full preparation after a
purchase, exists exactly when
\[
\DT>2k,\qquad \tfrac12<\tau\le M_{\min\{\lambda_R(\DT),\lb\}},
\]
where \(\lambda_R(\DT)\) is the unique root in \((0,1)\) of \eqref{eq:research-cond}. Its research is
\(\hat\lambda=\min\{\lambda_R(\DT),\lb\}\). Class \(H\) buys strictly and fully, class \(L\) uses Lemma A.15, and entry
is \(e(1)=1\), \(e(0)=e(-1)=\rho\). The posteriors are those of that lemma, prices are
\(P(x)=t_0+e(x)[t_L-t_0+\DT\mu(x)]\), price beliefs are \(\bar\mu(P(x))=\mu(x)\), and bids are truthful. Equal
nonpurchase posteriors have equal entry and equal prices.
\end{paperresult}

\paragraph{Proof, step by step.}
\begin{enumerate}
\item \emph{A strict full purchase.} Positive research forces \(\sigma_H(1)=1\) and \(U_H(1)>0\) (Lemma A.14). Entry
after a sale and after no order is \(\rho\).
\item \emph{The value of a success.} With full preparation after the purchase, \(U_H(1)=\DT m_{\hat\lambda}-k\).
Lemma A.15 gives class \(L\)'s best payoff. So \(\Psi=\tfrac12[\DT m_{\hat\lambda}-k+(\rho\DT m_{\hat\lambda}-k)^+]\).
\item \emph{The research level.} Section~\ref{sec:research-value} turns (E7) into \(\hat\lambda=\min\{\lambda_R,\lb\}\),
and the root needs \(\DT>2k\).
\item \emph{The purchase is strict at the capped choice.} At \(\hat\lambda\) the residual is nonnegative, so
\(\DT m-k+(\rho\DT m-k)^+\ge2a\hat\lambda>0\). If \(\DT m-k\le0\), then also \(\rho\DT m-k<0\), the positive part is
zero, and the left side is at most zero, a contradiction. So \(\DT m_{\hat\lambda}>k\).
\item \emph{Recruitment.} A purchase recruits exactly when \(M_{\hat\lambda}\ge\tau\). At equality the default rule
makes the challenger prepare; under the other rule, the selection does. With \(\tau>1/2\) from the setting, this is the
displayed condition.
\item \emph{The converse.} Under the displayed condition, build the tuple: research \(\hat\lambda\), \(\sigma_H=\delta_1\),
the sale rate of Lemma A.15, its posteriors, entry \((\rho,\rho,1)\), the price formula, beliefs equal to posteriors,
and truthful bids. Check each condition.
\begin{itemize}
\item (E1): the posteriors come from Bayes' rule.
\item (E2) and (E4): entry never falls when the posterior rises. For \(\mu(x)>\mu(y)\), the increment identity
\eqref{eq:beliefs-increment} gives \(P(x)-P(y)=e(x)\DT[\mu(x)-\mu(y)]+[e(x)-e(y)][t_L-t_0+\DT\mu(y)]>0\), because
\(t_L\ge t_0\) and \(e(x)\ge\rho>0\). A sale and no order can share a posterior; they then share entry \(\rho\) and a
price, and the belief at that price is their common posterior (Section~\ref{sec:beliefs-constructed}).
\item (E3): \(\mu(1)=M_{\hat\lambda}\ge\tau\), and the other posteriors are at most \(1/2<\tau\).
\item (E5): the purchase is strict for class \(H\) (step 4); class \(L\) follows Lemma A.15.
\item (E6): truthful bids.
\item (E7): the research equality, or the cap inequality, gives optimality of every joint deviation by Lemma A.12.
\end{itemize}
These conditions also fix every component, so the positive-research tuple is unique.
\end{enumerate}

\begin{checkyourself}\label{q:research-zero}
Take zero research, \(\hat\lambda=0\). Show that every posterior is \(1/2\), whatever the orders after a hypothetical
success, and that entry is \(\rho\). Find \(U_H(1)\) and \(U_L(-1)\), and show that zero research is an equilibrium
exactly when \(\rho\DT\le2k\). Why must the orders after a hypothetical success be optimal?
\end{checkyourself}

\subsection{Part (b): research and entry rise with the spread}

The appendix adds detail to part (b).

\begin{paperresult}{Proposition 4(b), full statement}
Research is continuous and strictly increasing in \(\DT\) below its cap, and constant at the cap. Total entry and class
entry are
\[
\Ent=\rho+(1-\rho)[\beta+(1/2-\beta)\hat\lambda],\qquad
e_H=\rho+(1-\rho)[\beta+(1-\beta)\hat\lambda],\qquad e_L=\rho+(1-\rho)\beta(1-\hat\lambda).
\]
Thus total entry and class-\(H\) entry rise, class-\(L\) entry falls, and the class-\(H\) share of prepared challengers
rises with the spread, strictly below the cap and weakly throughout the existence set. The prepared-value distribution
\((e_HG_H+e_LG_L)/(e_H+e_L)\) improves in first-order stochastic dominance; its mean rises strictly when research rises
strictly. For almost every fixed incumbent value \(x\), the high-class acquisition probability is \(e_H/2\) for
\(x<p\) and \(e_H[1-G_H(x)]/2\) for \(x\ge p\). It rises strictly with the spread below the cap wherever class \(H\) can
win.
\end{paperresult}

The appendix statement also gives the slopes of research, the kink at the sale boundary and the cap spread.
Section~\ref{sec:research-sale} derives them.

\paragraph{Step 1: research rises with the spread.} Write \(\mathcal R(\lambda;\DT)\) to show the spread. At a fixed
\(\lambda\in(0,1)\), \(m_\lambda>0\), so \(\DT m_\lambda\) rises strictly with \(\DT\), and the positive part never
falls. So \(\mathcal R\) rises strictly with \(\DT\). Take \(2k<\DT<\DT'\):
\begin{align*}
\mathcal R(\lambda_R(\DT);\DT')&>\mathcal R(\lambda_R(\DT);\DT) \why{\(\mathcal R\) rises strictly in \(\DT\)}\\
&=0 \why{definition of \(\lambda_R(\DT)\)}
\end{align*}
\(\mathcal R(\cdot;\DT')\) is strictly decreasing and zero at \(\lambda_R(\DT')\). It is positive at \(\lambda_R(\DT)\),
so \(\lambda_R(\DT)<\lambda_R(\DT')\).

\paragraph{Step 2: research moves continuously.} The paper states this; the guide gives the argument. Fix \(\DT>2k\)
and a small \(\varepsilon>0\). Strict decrease gives
\(\mathcal R(\lambda_R-\varepsilon;\DT)>0>\mathcal R(\lambda_R+\varepsilon;\DT)\). At each fixed \(\lambda\),
\(\mathcal R\) is continuous in \(\DT\). So both signs survive for every \(\DT'\) near enough to \(\DT\). Then the root
\(\lambda_R(\DT')\) lies between \(\lambda_R-\varepsilon\) and \(\lambda_R+\varepsilon\). Cutting at the cap,
\(\min\{\lambda_R,\lb\}\), keeps continuity and makes research constant at the cap.

\paragraph{Step 3: entry.} Only a purchase raises entry above \(\rho\). First find the probability of a purchase.

\begin{derivation}[entry with research]
\begin{align*}
\Pr(X=1)&=\tfrac12 f_H(1)+\tfrac12 f_L(1) \why{total probability over the classes}\\
&=\tfrac12(b+\hat\lambda)+\tfrac12 b=b+\tfrac12\hat\lambda \why{\(f_H(1)=b+\hat\lambda\), \(f_L(1)=b\)}\\
&=(1-\hat\lambda)\beta+\tfrac12\hat\lambda=\beta+(\tfrac12-\beta)\hat\lambda \why{\(b=(1-\hat\lambda)\beta\); collect \(\hat\lambda\)}\\[4pt]
\Ent&=1\cdot\Pr(X=1)+\rho\,[1-\Pr(X=1)] \why{entry \(1\) after a purchase, \(\rho\) otherwise}\\
&=\rho+(1-\rho)\,[\beta+(\tfrac12-\beta)\hat\lambda] \why{rearrange}
\end{align*}
In class \(H\), a purchase has probability \(b+\hat\lambda=\beta+(1-\beta)\hat\lambda\); in class \(L\), it has
probability \(b=\beta(1-\hat\lambda)\). The same steps give
\[
e_H=\rho+(1-\rho)[\beta+(1-\beta)\hat\lambda],\qquad e_L=\rho+(1-\rho)\beta(1-\hat\lambda).
\]
\end{derivation}

This is \eqref{eq:ccc-entry} at \(\pi=1/2\), with presence \(\hat\lambda\). The low class's sale rate does not enter.
It moves probability between a sale and no order, and both have entry \(\rho\).

The slopes in research are
\[
\frac{d\Ent}{d\hat\lambda}=(1-\rho)\bigl(\tfrac12-\beta\bigr)>0,\qquad
\frac{de_H}{d\hat\lambda}=(1-\rho)(1-\beta)>0,\qquad
\frac{de_L}{d\hat\lambda}=-(1-\rho)\beta<0 .
\]
The first is positive because \(\beta<1/2\). Combine with Steps 1 and 2: entry and \(e_H\) rise with the spread and
\(e_L\) falls, strictly below the cap and weakly over the whole existence set.

\begin{intuition}[one unit of research, two changes in purchases]
One more unit of research adds an informed purchase whenever the class is \(H\), which happens half the time. It also
removes a liquidity purchase with probability \(\beta\), in both classes. Since \(\beta<1/2\), the first change is
larger, and purchases become more frequent. Class \(L\) reaches a purchase only through a liquidity trader, so its
entry falls. The challengers who prepare are more often of class \(H\).
\end{intuition}

\paragraph{Step 4: the share of class \(H\).} Write \(S=e_H+e_L\). By the quotient rule,
\begin{align*}
\frac{d}{d\hat\lambda}\frac{e_H}{S}
&=\frac{e_H'\,S-e_H\,(e_H'+e_L')}{S^2}=\frac{e_H'\,e_L-e_H\,e_L'}{S^2} \why{quotient rule; \(e_H'e_H\) cancels}\\
&=\frac{(1-\rho)\bigl[(1-\beta)e_L+\beta e_H\bigr]}{S^2} \why{\(e_H'=(1-\rho)(1-\beta)\), \(e_L'=-(1-\rho)\beta\)}\\
&=\frac{(1-\rho)\bigl[\rho+(1-\rho)\beta\bigr]}{S^2}>0 \why{expand; the \(\hat\lambda\) terms cancel}
\end{align*}
In the last step, \((1-\beta)e_L+\beta e_H=\rho+(1-\rho)[(1-\beta)\beta(1-\hat\lambda)+\beta^2+\beta(1-\beta)\hat\lambda]
=\rho+(1-\rho)\beta\).

\paragraph{Step 5: the values of prepared challengers.} Preparation depends only on the order and the cost, so given
the class it is independent of the challenger's value. A prepared challenger is of class \(j\) with probability
\(e_j/S\). So its value has the distribution \(sG_H+(1-s)G_L\), with \(s=e_H/S\). The class ranking gives
\(G_H\le G_L\), so a higher \(s\) lowers this mixture at every point: a first-order improvement. Its mean
\(s\,\E_H\theta+(1-s)\,\E_L\theta\) rises at the rate \(\E_H\theta-\E_L\theta=\DT+\gamma>0\) per unit of \(s\)
(Lemma A.5).

\paragraph{Step 6: acquisitions at a fixed incumbent value.} The same independence gives
\(\Pr(H\text{ acquires}\mid R=x)=\tfrac12\,e_H\,\Pr(\theta>x\mid H)\) for \(x\ge p\), and \(\tfrac12 e_H\) for \(x<p\),
where the incumbent cannot meet the reserve. The win factor does not depend on research, so the change has the sign of
the change in \(e_H\) wherever class \(H\) can win.

\begin{checkyourself}\label{q:research-entry}
Entry in class \(L\) falls with research, although a purchase recruits. Explain why with the formula for \(e_L\). What
would total entry do as research rises if \(\beta\) were \(1/2\), and why does the model exclude this value?
\end{checkyourself}

\subsection{Part (c): ownership}

\begin{paperresult}{Proposition 4(c), full statement}
High-class ownership is \(\Own_H(F)=e_H(F)W_H(F)/2\). For two recruiting economies on a path,
\[
\Own_H(F_{r'})\ge\Own_H(F_r)\quad\Longleftrightarrow\quad\frac{W_H(F_{r'})}{W_H(F_r)}\ge\frac{e_H(F_r)}{e_H(F_{r'})},
\]
with the corresponding strict equivalence. These win probabilities are positive. If class \(H\) outbids every
incumbent along the recruiting part of the path, ownership rises weakly with strength, and strictly for a strict
increase in the spread from a research intensity below the cap. Decisive advantage (Lemma A.6),
\(g_H(F)>c_H\Rightarrow W_H(F)=1\) throughout the family, supplies this outbidding condition at every recruiting
economy.
\end{paperresult}

\paragraph{Proof, step by step.}
\begin{enumerate}
\item \emph{The formula.} Given the class, preparation is independent of the values. So
\(\Own_H=\Pr(H)\cdot\Pr(\text{preparation}\mid H)\cdot W_H(F)=\tfrac12\,e_H\,W_H(F)\).
\item \emph{Positive win probabilities.} At a recruiting economy, \(\tau\le M_{\hat\lambda}\), so
\(c_H\le B(M_{\hat\lambda})\). Then
\begin{align*}
c_H\le B(M_{\hat\lambda})&=g_H-(1-M_{\hat\lambda})\gamma \why{\(B(\mu)=g_L+\mu\gamma\); rearrange}\\
&<g_H \why{\(M_{\hat\lambda}<1\), \(\gamma>0\)}
\end{align*}
So \(g_H>c_H>0\). A positive gross profit means a positive chance to win (Chapter~\ref{ch:ccc}, Step 4 of part (ii)), so
\(W_H(F)>0\).
\item \emph{The equivalence.} \(\Own_H(F_{r'})\ge\Own_H(F_r)\) means \(e_H(F_{r'})W_H(F_{r'})\ge e_H(F_r)W_H(F_r)\). Divide
by the positive number \(W_H(F_r)\,e_H(F_{r'})\). The strict version has the same steps.
\item \emph{Outbidding.} If \(W_H=1\) everywhere, then \(\Own_H=e_H/2\): half the entry in class \(H\). Along a path the
spread never falls (Lemma A.5), so research never falls, and by part (b) \(e_H\) never falls. Where research is below the
cap and the spread rises strictly, research and \(e_H\) rise strictly.
\item \emph{Decisive advantage.} Step 2 gives \(g_H>c_H\) at every recruiting economy. In a family with decisive
advantage this gives \(W_H=1\) (Section~\ref{sec:ccc-decisive}).
\end{enumerate}

\paragraph{Ownership at the same incumbent.} The appendix compares two equilibria at one incumbent. Where recruiting
research coexists with zero research, \(\rho\DT\le2k\), zero research has entry \(\rho\) after every order, so
\(\Own_H^{\rm Z}=\rho W_H(F)/2\). Divide: the win probability cancels, and
\[
\frac{\Own_H^{\rm R}}{\Own_H^{\rm Z}}=\frac{e_H}{\rho}>1 .
\]
This ratio rises strictly with the spread below the cap and stays constant at the cap.

\begin{watchout}
``High-class ownership equals half the entry in class \(H\)'' needs class \(H\) to outbid the incumbent,
\(W_H(F)=1\). Without it, ownership rises only under the ratio condition of the full statement, because a stronger
incumbent can lower class \(H\)'s chance to win. Decisive advantage is one way to obtain the outbidding.
\end{watchout}

\subsection{Part (d): one interval of strengths}

A path holds every other primitive fixed and strengthens the incumbent as \(r\) rises. Along it the spread never
falls, and, where \(\tau>1/2\), the threshold never falls (Lemma A.5, Section~\ref{sec:trading-path}). Call the
spread at which research reaches its cap the \emph{cap spread}.

\begin{paperresult}{Proposition 4(d), full statement}
Let \(I=[r_-,r_+]\) be compact and \(\DT,\tau\) continuous along the path. The capped recruiting set is
\[
\bigl\{r\in I:\ \DT(r)m_{\lb}-k+(\rho\DT(r)m_{\lb}-k)^+\ge2a\lb,\ \ \tau(r)\le M_{\lb}\bigr\}.
\]
It is empty, a singleton, or a nondegenerate closed interval. If nonempty, it ends at \(\max\{r\in I:\tau(r)\le M_{\lb}\}\).

Suppose the path first reaches the cap spread at \(r_{\lb}>r_-\), with a strictly smaller spread at \(r_-\). On
\([r_-,r_{\lb}]\), suppose \(\tau<1\), and define \(\varphi_a^*=J-k+(\rho J-k)^+-2a\lambda_E(\tau)\). If \(\varphi_a^*\)
is strictly increasing on \([r_-,r_{\lb}]\) and \(\varphi_a^*(r_-)<0<\varphi_a^*(r_{\lb})\), the full recruiting set is
\[
\bigl[r_R,\ \max\{r\in I:\tau(r)\le M_{\lb}\}\bigr],
\]
where \(r_R\) is the unique zero of \(\varphi_a^*\) before \(r_{\lb}\). Research is below its cap up to \(r_{\lb}\) and at
the cap thereafter. A threshold plateau at \(M_{\lb}\) belongs to the recruiting set through its last point. On a path
with continuously differentiable \(\DT,\tau\) on \([r_-,r_{\lb}]\), the cheap-research condition
\[
\text{(CR)}\qquad [1+\rho\1\{\rho J>k\}]\,J'-\frac{2a\beta\,\tau'}{[\beta+(1-2\beta)(1-\tau)]^2}>0
\quad\text{almost everywhere on }[r_-,r_{\lb}]
\]
makes \(\varphi_a^*\) strictly increasing there.
\end{paperresult}

\paragraph{Step 1: research at the cap.} At the cap, research is optimal exactly when the residual at \(\lb\) is
nonnegative:
\[
\DT m_{\lb}-k+(\rho\DT m_{\lb}-k)^+\ge2a\lb .
\]
The left side rises strictly and continuously with \(\DT\), so equality picks out one spread, the cap spread. At or
above the cap spread, the purchase pays strictly, by the argument of step 4 of part (a). The purchase posterior is
\(M_{\lb}\), so it recruits exactly when \(\tau\le M_{\lb}\). Bayes' rule, the price formula and truthful bids give the
other conditions.

\paragraph{Step 2: the capped recruiting set.} Take two kinds of sets of strengths.
\begin{itemize}
\item \(\{r:\DT(r)\ \text{is at least the cap spread}\}\). The spread never falls, so this set contains every strength
above each of its members.
\item \(\{r:\tau(r)\le M_{\lb}\}\). The threshold never falls, so this set contains every strength below each of its
members.
\end{itemize}
Continuity makes both sets closed. So the first is an interval that runs to \(r_+\), and the second an interval that
starts at \(r_-\). Their intersection is empty, one point, or a closed interval. If it is not empty, it ends where the
second set ends, at \(\max\{r:\tau(r)\le M_{\lb}\}\). A stretch with \(\tau=M_{\lb}\) lies in the second set, so it
belongs to the interval through its last point.

\paragraph{Step 3: the margin decides below the cap.} For \(r\in[r_-,r_{\lb}]\), the derivation after
\eqref{eq:research-margin} gives: a purchase at the chosen research recruits exactly when \(\varphi_a^*\ge0\) and
\(\lambda_E(\tau)\le\lb\). The margin is continuous, strictly increasing, negative at \(r_-\) and positive at
\(r_{\lb}\). So it has exactly one zero \(r_R\), and it is negative before \(r_R\) and positive after.

\paragraph{Step 4: the cap condition holds below \(r_{\lb}\).} At \(r_{\lb}\), continuity makes the spread equal to the
cap spread, so \(\lambda_R=\lb\). The positive margin there gives \(\mathcal R(\lambda_E)>0\), so
\(\lambda_E<\lambda_R=\lb\), that is \(\tau(r_{\lb})<M_{\lb}\). The threshold never falls, so \(\tau(r)<M_{\lb}\) on all
of \([r_-,r_{\lb}]\), and \(\lambda_E(\tau)<\lb\) there. On this stretch, recruitment is therefore the same as
\(\varphi_a^*\ge0\), that is \(r\ge r_R\). Before \(r_{\lb}\) the spread is below the cap spread, so research is below
the cap.

\paragraph{Step 5: beyond \(r_{\lb}\).} The spread never falls, so research stays at the cap from \(r_{\lb}\) on. Step 2
gives recruitment exactly up to \(\max\{r:\tau(r)\le M_{\lb}\}\). Together, the recruiting set is
\([r_R,\max\{r:\tau(r)\le M_{\lb}\}]\).

\paragraph{Step 6: the cheap-research condition.} Away from \(\rho J=k\), differentiate \eqref{eq:research-margin}
along the path:
\begin{align*}
\frac{d\varphi_a^*}{dr}&=J'+\rho J'\,\1\{\rho J>k\}-2a\,\frac{d\lambda_E}{d\tau}\,\tau' \why{chain rule; \((\rho J-k)^+\) has slope \(\rho J'\) where \(\rho J>k\)}\\
&=[1+\rho\1\{\rho J>k\}]\,J'-\frac{2a\beta\,\tau'}{[\beta+(1-2\beta)(1-\tau)]^2} \why{slope of \(\lambda_E\), Section~\ref{sec:research-posteriors}}
\end{align*}
On a stretch where \(\rho J=k\) holds, \(J\) is constant, so \(J'=0\) almost everywhere there, and the same expression
is still the derivative. A function whose derivative is positive almost everywhere rises strictly over every stretch,
because it is the integral of its derivative. So (CR) makes the margin rise strictly.

\begin{watchout}
Part (d) compares economies along a path. In each economy the equilibrium is solved again: the spread, the threshold,
research and the prices all change. Inside one economy, a deviation in research faces fixed prices and entry.
\end{watchout}

% =====================================================================
\section{The low class's sale}\label{sec:research-sale}

Class \(L\) sells after bad news when a sale at the entry floor pays. The paper explains why this leaves the mechanism
in place.
\begin{itemize}
\item \emph{A sale never recruits.} Its posterior is at most \(1/2<\tau\), so its entry is \(\rho\).
\item \emph{While class \(L\) mixes, its sale breaks even.} Its best payoff is zero, as when it abstains. So a success
has the same value, and \eqref{eq:research-cond} has no sale term.
\item \emph{Once a sure sale pays,} its profit \(\rho\DT m_\lambda-k\) is a second return to research. If research is
still below its cap at that point, the research schedule kinks upward there.
\end{itemize}

\paragraph{The two slopes of research.} On each side of the sale boundary, \eqref{eq:research-cond} is smooth.
\begin{itemize}
\item While class \(L\) abstains or mixes: \(2a\lambda=\DT m_\lambda-k\).
\item After a sure sale pays strictly: \(2a\lambda=\DT m_\lambda-k+\rho\DT m_\lambda-k=(1+\rho)\DT m_\lambda-2k\).
\end{itemize}
Differentiate each with respect to \(\DT\), with \(\lambda=\lambda_R(\DT)\) and \(m'=dm_\lambda/d\lambda<0\).

\begin{derivation}[the slopes of research in the spread]
First piece:
\begin{align*}
2a\,\lambda_R'&=m_{\lambda_R}+\DT\,m'_{\lambda_R}\,\lambda_R' \why{product rule and chain rule}\\
\lambda_R'\,\bigl(2a-\DT m'_{\lambda_R}\bigr)&=m_{\lambda_R} \why{collect the terms in \(\lambda_R'\)}\\
\lambda_R'&=\frac{m_{\lambda_R}}{2a-\DT m'_{\lambda_R}} \why{divide; the bracket exceeds \(2a>0\)}
\end{align*}
Second piece, with the factor \(1+\rho\) on \(\DT m_\lambda\):
\[
\lambda_R'=\frac{(1+\rho)\,m_{\lambda_R}}{2a-(1+\rho)\DT m'_{\lambda_R}} .
\]
\end{derivation}

Both slopes are positive, because \(m>0\) and \(m'<0\). This confirms part (b) on each piece.

\paragraph{Where the pieces meet.} The sale boundary is \(\rho\DT m_{\lambda_R}=k\). Combine it with the first piece:
\begin{align*}
\DT m_{\lambda_R}&=k/\rho \why{sale boundary}\\
2a\lambda_R&=k/\rho-k=\frac{(1-\rho)k}{\rho} \why{first piece}\\
\lambda_R&=\frac{(1-\rho)k}{2a\rho} \why{divide by \(2a\)}
\end{align*}
At this point, write \(A=2a\) and \(G=-\DT m'_{\lambda_R}>0\). Subtract the first slope from the second:
\begin{align*}
\frac{(1+\rho)m}{A+(1+\rho)G}-\frac{m}{A+G}
&=\frac{m\,[(1+\rho)(A+G)-A-(1+\rho)G]}{(A+G)[A+(1+\rho)G]} \why{common denominator}\\
&=\frac{m\,\rho A}{(A+G)[A+(1+\rho)G]}=\frac{2a\rho\,m_{\lambda_R}}{(2a-\DT m'_{\lambda_R})[2a-(1+\rho)\DT m'_{\lambda_R}]}>0
\why{expand; \(A+G\) terms cancel}
\end{align*}
So, if \((1-\rho)k/(2a\rho)<\lb\), the slope of research jumps up at this intensity. Research itself and entry stay
continuous, through the first sale, through the onset of a strictly profitable sure sale, and through the cap.

\paragraph{The cap spread.} Set \(\lambda=\lb\) in the research condition on each side of the sale boundary.
\begin{itemize}
\item Sale term zero: \(2a\lb=\DT m_{\lb}-k\) gives \(\DT=(2a\lb+k)/m_{\lb}\).
\item Sure sale pays: \(2a\lb=(1+\rho)\DT m_{\lb}-2k\) gives \(\DT=(2a\lb+2k)/[(1+\rho)m_{\lb}]\).
\end{itemize}
The first case applies when the sale term is zero at its own spread: \(\rho(2a\lb+k)\le k\), that is
\(\lb\le(1-\rho)k/(2a\rho)\). Both pieces agree at the boundary. The cap spread exceeds \(2k\): at \(\DT=2k\), the
residual at the cap is \(k(2m_{\lb}-1)-2a\lb<0\), because \(m_{\lb}<1/2\) and the sale term is zero there.

\begin{intuition}[a second return, if research still has room]
While the low class mixes, the market reads a sale as just enough news to make it break even, so a sale adds no value
to a success. A sure sale that pays is different: a success then earns something in both classes. Research responds
faster to the spread from that point on. If research already sits at its cap, nothing changes. In the example,
research reaches the cap before the low class starts to sell, so the sale's profit leaves research unchanged.
\end{intuition}

\begin{checkyourself}\label{q:research-kink}
Suppose \((1-\rho)k/(2a\rho)\ge\lb\). Show that research reaches the cap at a spread where a sure sale does not pay
strictly. Explain in words why the low class's sale then never changes research.
\end{checkyourself}

% =====================================================================
\section{The race between the spread and the threshold}\label{sec:research-race}

Section~\ref{sec:trading-race} described a race along a path. A stronger incumbent never narrows the spread, and,
while \(g_L\le c_H\), it never lowers the threshold. With research, the race reads as follows.
\begin{itemize}
\item \emph{The spread.} A wider spread raises what a recruiting purchase earns, \(J=\DT(1-\tau)\) at a fixed
threshold, and so the value of a success.
\item \emph{The threshold.} A higher threshold raises the belief that a purchase must reach. So it raises the
research needed, \(\lambda_E(\tau)\), and the marginal cost \(a\lambda_E\) of that research. It also lowers \(J\) at a
fixed spread.
\end{itemize}
The exact recruitment margin \eqref{eq:research-margin} collects both forces in one number. Part (d) asks that it
rise strictly up to \(r_{\lb}\). Then, once below the cap, the race tips toward recruitment, and the equilibrium holds
on one interval of strengths. It lasts through the last strength at which the threshold does not exceed the top
posterior at the cap.

\paragraph{The cheap-research condition in words.} Halve (CR). Its first term, \(\tfrac12[1+\rho\1\{\rho J>k\}]J'\), is
the rise in the value of a success at the recruitment intensity. Its second term is the rise in the marginal cost,
\(a\,\lambda_E'(\tau)\,\tau'\). So (CR) says: the increase in a recruiting purchase's return exceeds the marginal cost
of the additional research needed to reach the threshold. Where \(\tau'>0\), it is an upper bound on the cost:
\[
a<\frac{[1+\rho\1\{\rho J>k\}]\,J'}{2\,\lambda_E'(\tau)\,\tau'} .
\]
Where \(\tau'=0\), it says \(J'>0\).

The paper adds two remarks on the margin. Before a sure sale pays, its derivative is \(J'-2a\lambda_E'(\tau)\tau'\);
once a sure sale pays, it is \((1+\rho)J'-2a\lambda_E'(\tau)\tau'\). The margin is continuous at \(\rho J=k\), and at a
crossing into \(\rho J>k\) its derivative jumps up by \(\rho J'\).

\paragraph{Extreme cases.} With research cost near zero, the condition holds wherever \(J\) rises. With a large cost,
it can fail even where \(J\) rises, because each rise in the threshold then needs costly research. Recall from
Section~\ref{sec:trading-path} that \(J\) itself can rise or fall along a path.

% =====================================================================
\section{The three research outcomes}\label{sec:research-outcomes}

Proposition~4 covers one equilibrium. The appendix lists every equilibrium under the default tie rules. Two more
outcomes appear. In \emph{zero research}, the investor does no research. In \emph{floor-entry research}, both classes
trade after a success, no order recruits, and entry is \(\rho\).

For \(\rho\DT\ge2k\), the paper defines \(\lambda_Q(\DT)\) as the unique root in \([0,1)\) of
\begin{equation}\label{eq:research-floor}
a\lambda_Q=\rho\DT\,m_{\lambda_Q}-k .
\end{equation}
It equals zero exactly when \(\rho\DT=2k\).

\begin{paperresult}{Lemma A.16 (The three research outcomes)}
Maintain the setting above. An indifferent investor may mix over its best orders and an indifferent high-cost
challenger prepares. The following outcomes exhaust the equilibria.
\begin{enumerate}[label=(\alph*),leftmargin=2em]
\item Zero research exists exactly when \(\rho\DT\le2k\). All posteriors are \(1/2\), entry is \(\rho\), and every order
has price \(t_0+\rho[t_L-t_0+\DT/2]\). If the inequality is strict, both classes abstain after a hypothetical success.
At equality, each class uses any mixture of its correct order and abstention; all these strategies give the same
zero-research outcome.
\item Recruiting research is exactly the outcome of Proposition~4(a).
\item Floor-entry research exists exactly when \(\rho\DT>2k\) and \(\tau>M_{\min\{\lambda_Q(\DT),\lb\}}\). Research is
\(\hat\lambda=\min\{\lambda_Q,\lb\}\), both classes trade strictly and fully after a success, the posteriors are
\((m_{\hat\lambda},1/2,M_{\hat\lambda})\), and entry is \(\rho\) after every order. Prices are
\(t_0+\rho[t_L-t_0+\DT\mu(x)]\), price beliefs are the corresponding posteriors, and bids are truthful.
\end{enumerate}
Each positive-research outcome has a unique tuple. The two positive-research outcomes coexist exactly when
\[
\rho\DT>2k,\qquad M_{\min\{\lambda_Q,\lb\}}<\tau\le M_{\min\{\lambda_R,\lb\}}.
\]
The uncapped roots satisfy \(\lambda_Q<\lambda_R\); their capped values are strictly ordered when \(\lambda_Q<\lb\) and
coincide when both roots reach the cap. Thus the coexistence band has positive width in \(\tau\) exactly when
\(\lambda_Q<\lb\). At least one of the three outcomes exists at every auction in the stated domain, including
\(\tau\ge1\).
\end{paperresult}

\paragraph{Proof, step by step.}
\begin{enumerate}
\item \emph{Zero research.} Check yourself~\ref{q:research-zero} works this case in full. At \(\hat\lambda=0\) every
order comes from a liquidity trader in both classes, so every posterior is \(1/2\) and entry is \(\rho\). The correct
orders earn \(\rho\DT/2-k\), so \(\Psi=(\rho\DT/2-k)^+\). Zero research is the best choice exactly when \(\Psi=0\), that
is \(\rho\DT\le2k\). At equality both correct orders tie with abstention. They have probability zero, so any mixture
leaves beliefs, prices and entry unchanged.
\item \emph{Every positive-research equilibrium.} Lemma A.14 gives a strict full purchase, \(\mu(1)=M_{\hat\lambda}\),
and entry \(\rho\) after the other orders. Two cases remain.
\begin{itemize}
\item If \(M_{\hat\lambda}\ge\tau\), the purchase recruits. This is recruiting research (part (b) of the lemma).
\item If \(M_{\hat\lambda}<\tau\), the purchase earns \(\rho\DT m_{\hat\lambda}-k>0\). A sale earns at least as much,
because \(\mu(-1)\ge m_{\hat\lambda}\). So class \(L\) sells strictly and fully, and both classes earn
\(\rho\DT m_{\hat\lambda}-k\). Then \(\Psi=\rho\DT m_{\hat\lambda}-k\), and Lemma A.12 gives
\(\hat\lambda=\min\{\lambda_Q,\lb\}\). This is floor-entry research.
\end{itemize}
\item \emph{The floor root.} The gap \(\rho\DT m_\lambda-k-a\lambda\) falls strictly in \(\lambda\), from
\(\rho\DT/2-k\ge0\) at \(\lambda=0\) to \(-k-a\) at \(\lambda=1\). So \(\lambda_Q\) exists and is unique. It is positive
exactly when \(\rho\DT>2k\). The gap rises strictly with \(\DT\), so \(\lambda_Q\) rises strictly and continuously with
the spread, by the argument of part (b).
\item \emph{The order of the roots.} Let \(\rho\DT>2k\). At \(\lambda_Q\), the floor equation gives
\(\rho\DT m_{\lambda_Q}-k=a\lambda_Q>0\), so the positive part in \(\mathcal R\) is active. Then
\begin{align*}
\mathcal R(\lambda_Q)&=\DT m_{\lambda_Q}-k+(\rho\DT m_{\lambda_Q}-k)-2(\rho\DT m_{\lambda_Q}-k) \why{\(2a\lambda_Q=2(\rho\DT m_{\lambda_Q}-k)\)}\\
&=(1-\rho)\,\DT\,m_{\lambda_Q}>0 \why{the \(k\) terms cancel}
\end{align*}
Since \(\mathcal R\) is strictly decreasing with root \(\lambda_R\), \(\lambda_Q<\lambda_R\).
\item \emph{Coexistence.} Floor-entry research needs \(\tau>M_{\min\{\lambda_Q,\lb\}}\); recruiting research needs
\(\tau\le M_{\min\{\lambda_R,\lb\}}\). Both hold on the displayed band. If both roots reach the cap, the two bounds
coincide, and the band is empty.
\item \emph{Some outcome always exists.} If \(\rho\DT\le2k\), zero research exists. Otherwise, either
\(\tau\le M_{\min\{\lambda_R,\lb\}}\) and recruiting research exists, or \(\tau>M_{\min\{\lambda_R,\lb\}}\ge
M_{\min\{\lambda_Q,\lb\}}\) and floor-entry research exists. This covers \(\tau\ge1\) too.
\end{enumerate}
The construction of floor-entry research checks (E1) to (E7) as in part (a): full orders give the three posteriors,
constant entry \(\rho\) makes the three prices distinct, and the strict threshold inequality gives preparation.

\begin{intuition}[two research levels that confirm themselves]
This reading is the guide's own. If the market expects recruiting research, a purchase brings full entry, so a success
is worth more and the investor researches more. Its purchase posterior then reaches the threshold. If the market
expects floor-entry research, a purchase brings only the floor, a success is worth less, and the lower research keeps
the purchase posterior below the threshold. On the band of Lemma A.16, both expectations are correct.
\end{intuition}

\begin{checkyourself}\label{q:research-floor}
In floor-entry research, show that the value of a success is \(\Psi=\rho\DT m_\lambda-k\), and derive
\eqref{eq:research-floor}. Why does \eqref{eq:research-floor} have the factor \(a\), and \eqref{eq:research-cond} the
factor \(2a\)?
\end{checkyourself}

\paragraph{Strong incumbents.} Against a strong enough incumbent, the equilibrium is unique under every tie rule.

\begin{paperresult}{Lemma A.17 (Strong incumbents)}
Maintain the setting above. If \(M_{\lb}\ne\tau\) and
\[
\text{(U)}\qquad \rho\DT m_{\lb}-k>a\lb ,
\]
then under all four tie rules the equilibrium is unique. Research is \(\lb\), class \(H\) buys and class \(L\) sells
strictly after a success, and the posteriors are \((m_{\lb},1/2,M_{\lb})\). Entry after a sale and no order is \(\rho\),
while \(\zeta(P(1))=\1\{M_{\lb}>\tau\}\). Competitive prices, their posterior beliefs and truthful bids complete the
tuple. With \(\bar b=(1-\lb)\beta\), total entry is \(\rho+(1-\rho)(\bar b+\lb/2)\) if \(M_{\lb}>\tau\) and \(\rho\)
otherwise.
\end{paperresult}

\paragraph{Proof, step by step.} Take any candidate.
\begin{enumerate}
\item Lemma A.13(a) puts every nonzero-order posterior in \([m_{\hat\lambda},M_{\hat\lambda}]\subset[m_{\lb},M_{\lb}]\),
whatever the orders, and entry is at least \(\rho\).
\item So \(U_H(1)\ge\rho\DT(1-M_{\lb})-k=\rho\DT m_{\lb}-k\) and \(U_L(-1)\ge\rho\DT m_{\lb}-k\). By (U), both exceed
\(a\lb\ge0\). The purchase and the sale are the strict best orders, under both investor rules.
\item Then \(\Psi=\tfrac12[U_H(1)+U_L(-1)]>a\lb\), and Lemma A.12 gives \(\hat\lambda=\lb\).
\item Full orders give the posteriors \((m_{\lb},1/2,M_{\lb})\) (Lemma A.1(d) at the prior \(1/2\)).
\item Since \(M_{\lb}\ne\tau\), preparation after a purchase is one or zero with no tie, under both challenger rules.
The purchase has probability \(\bar b+\lb/2\), and it recruits exactly when \(M_{\lb}>\tau\).
\end{enumerate}
In words, (U) says that a full order earns more than the marginal cost of research at the cap, even at the entry
floor and at the least favourable posterior. Then nothing about the other players' behaviour can stop full orders and
full research.

% =====================================================================
\section{Corollary 3: research in the example}\label{sec:research-example}

The example has point-mass classes \(h\) and \(\ell\), a reserve \(p\), and an incumbent uniform on \([0,r]\), with
\(p<\ell<r<h\) on the path \([r_0,r_2]\). Chapter~\ref{ch:example} derives its formulas; on the path
\[
\DT(r)=\frac{(r-\ell)^2}{2r},\qquad g_L(r)=\frac{\ell^2-p^2}{2r},\qquad g_H(r)=h-\frac r2-\frac{p^2}{2r}.
\]
The spread and the threshold rise strictly on the path, with \(1/2<\tau<1\). The research cap is
\(\lb=u\beta/(1+u\beta)\) with \(1+u=\eu\), so \(m_{\lb}=1/(1+\eu)\) and \(M_{\lb}=\eu/(1+\eu)\)
(Section~\ref{sec:research-posteriors}). Chapter~\ref{app:numbers} lists the values of \(a\), \(\lb\) and the marked
strengths.

\subsection{The marked strengths}

The paper writes \(\lambda_E(r)=\lambda_E(\tau(r))\) and uses two margins along the path:
\[
\varphi_a(r)=\DT(r)[1-\tau(r)]-k-2a\lambda_E(r),\qquad \varphi_\rho(r)=\rho\DT(r)[1-\tau(r)]-k-a\lambda_E(r).
\]
The first is the exact recruitment margin without its sale term. The second, the \emph{floor margin}, is the gap of
\eqref{eq:research-floor} at \(\lambda=\lambda_E(r)\): the return to a full order at the entry floor, at the research
that lets a purchase just reach the threshold, net of the trading cost and the marginal cost of that research. Each
strength below is the unique root on the path of its equation.
\begin{itemize}
\item \(r_R\): \(\varphi_a(r_R)=0\). A researched purchase first recruits here. A purchase that just reaches the
threshold also pays for its research: \(J=\DT(1-\tau)=k+2a\lambda_E\). The factor \(2\) appears because the purchase
comes after half of the successes.
\item \(r_{\lb}\): \(\DT(r_{\lb})=(2a\lb+k)(1+\eu)\), the cap spread with no sale term. Research reaches its cap here.
\item \(r_N\): \(\rho\DT(r_N)=2k\). Against a quiet market, a unit order at the entry floor earns exactly its cost.
Zero research stops at \(r_N\), and in recruiting research the low class starts to sell.
\item \(r_S\): \(\rho\DT(r_S)\,m_{\lb}=k\). A sure sale at the cap just pays. The low class sells for sure from here.
\item \(r_Q\): \(\varphi_\rho(r_Q)=0\). The research of floor-entry research reaches \(\lambda_E\), so its purchase
posterior reaches the threshold. Floor-entry research below the cap ends here.
\item \(r_U\): \(\rho\DT(r_U)\,m_{\lb}-k=a\lb\), the edge of condition (U). The floor root reaches the cap here.
\item \(r_C\): \(\tau(r_C)=M_{\lb}=\eu/(1+\eu)\), the entry ceiling. Research at the cap lets a purchase just reach
the threshold.
\end{itemize}
The cap equals the example's presence probability, so \(r_N\), \(r_S\) and \(r_C\) are the strengths of the same names
in Chapter~\ref{ch:coexistence}. The paper shows the order
\[
r_0<r_R<r_{\lb}<r_N<r_S<r_Q<r_U<r_1<r_C<r_2 .
\]
The three intermediate strengths of the paper's table lie between \(r_R\) and \(r_{\lb}\). There, research satisfies
\(\lambda_E<\lambda_R<\lb\): it is positive and above the recruitment intensity, and it is below the cap.

\subsection{The statement}

\begin{paperresult}{Corollary 3 (Research equilibria in the example)}
In the example at the prior \(1/2\), with the paper's research cost \(a\) and cap \(\lb\), let an indifferent investor
mix over its best orders and an indifferent high-cost challenger prepare. On the path from \(r_0\) to \(r_2\):
\begin{enumerate}[label=(\roman*),leftmargin=2em]
\item zero research, with entry \(\rho\), is an equilibrium exactly on \([r_0,r_N]\);
\item recruiting research is an equilibrium exactly on \([r_R,r_C]\); research is below the cap on \([r_R,r_{\lb})\) and
at the cap on \([r_{\lb},r_C]\), and entry rises from \(\rho+(1-\rho)[\beta+(1/2-\beta)\lambda_E(r_R)]\) at \(r_R\) to
\(\rho+(1-\rho)[\beta+(1/2-\beta)\lb]\) at \(r_{\lb}\) and keeps that value up to \(r_C\);
\item floor-entry research, in which both classes trade after a success, no order recruits and entry is \(\rho\), is an
equilibrium exactly on \((r_N,r_Q)\), below the cap, and on \((r_C,r_2]\), at the cap;
\item there is no other equilibrium.
\end{enumerate}
The unique equilibrium thus has zero research at \(r_0\), recruiting research at the cap at \(r_1\) and floor-entry
research at the cap at \(r_2\); at \(r_1\) and \(r_2\) it is the equilibrium of Proposition~2 with \(\lambda=\lb\).
\end{paperresult}

The main text prints values for \(a\), \(\lb\), the strengths and the two entry levels; the box writes them in symbols.
The appendix version adds the following.
\begin{itemize}
\item The equilibrium set is: zero research alone on \([r_0,r_R)\); zero research and recruiting research on
\([r_R,r_N]\); recruiting research at the cap and floor-entry research on \((r_N,r_Q)\); recruiting research at the cap
alone on \([r_Q,r_C]\); floor-entry research at the cap alone on \((r_C,r_2]\).
\item In recruiting research, class \(L\) abstains through \(r_N\), mixes on \((r_N,r_S)\), and sells with probability
one on \([r_S,r_C]\). Floor-entry research has full orders and entry \(\rho\) throughout.
\item At \(r_N\), zero research has all the hypothetical-success order mixtures of Lemma A.16(a).
\item At each of the three intermediate strengths, exactly two equilibria coexist: zero research, and recruiting
research with \(\lambda_E<\lambda_R<\lb\), in which class \(H\) buys and class \(L\) abstains after a success. At
\(r_1\) and \(r_2\), the unique equilibrium has full orders and research at the cap, with entry
\(\rho+(1-\rho)[\beta+(1/2-\beta)\lb]\) at \(r_1\) and \(\rho\) at \(r_2\).
\end{itemize}

A plain reading. Against a weak incumbent, no informed order pays, and the investor does no research. As the
incumbent grows stronger, the spread widens. From \(r_R\), a purchase can pay for the research that makes it recruit,
and recruiting research appears beside zero research. Its research and entry rise until the cap. From \(r_N\), an
informed order pays even at the floor, so zero research disappears and floor-entry research appears. From \(r_Q\),
recruiting research is the only equilibrium. Beyond \(r_C\), even research at the cap leaves the purchase posterior
below the threshold, and entry returns to the floor.

\subsection{The proof in outline}

\paragraph{Step 1: recruiting research at the cap.} Substitute \(\lb\) in \(m_\lambda\): \(m_{\lb}=1/(2+u)=1/(1+\eu)\).
In the example the ratio \((1-\rho)k/(2a\rho)\) exceeds one, so it exceeds the cap, and the cap spread is
\((2a\lb+k)(1+\eu)\) (Section~\ref{sec:research-sale}). Part (d), Step 2, gives recruiting research at the cap exactly on
\([r_{\lb},r_C]\), because \(r_{\lb}<r_N<r_1<r_C\). For the same reason the sale term in \eqref{eq:research-cond}
vanishes at every uncapped root, and the explicit root of Section~\ref{sec:research-value} applies.

\paragraph{Step 2: recruiting research below the cap.} On \([r_0,r_N]\), \(\rho\DT\le2k\) and \(1-\tau<1/2\), so
\(\rho J<\rho\DT/2\le k\). The sale term is zero, and \(\varphi_a^*=\varphi_a\). The paper bounds the derivatives of
\(\DT\) and \(\tau\) on a stretch of the path that contains \([r_0,r_N]\), and shows that \(\varphi_a\) rises there. The
same bound holds for every research cost below a level that the paper computes. Also \(\varphi_a(r_0)<0<\varphi_a(r_{\lb})\).
Part (d) then gives recruiting research on \([r_R,r_C]\), below the cap until \(r_{\lb}\). Beyond that stretch,
\(\varphi_a\) stays positive, because \(J\) rises on the whole path and \(\lambda_E\) is at most \(\lambda_E(r_2)\). So
\(r_R\) is the only zero of \(\varphi_a\) on the path.

\paragraph{Step 3: zero research.} Lemma A.16(a) gives zero research exactly where \(\rho\DT\le2k\), that is on
\([r_0,r_N]\).

\paragraph{Step 4: floor-entry research (Lemma A.18).} The paper shows that the floor margin \(\varphi_\rho\) rises
strictly on the path. It bounds the derivative of \(\varphi_\rho\) from below on short pieces of the path, with exact
arithmetic. Then:
\begin{itemize}
\item At \(r_N\), \(\rho\DT=2k\), so \(\varphi_\rho(r_N)=k(1-2\tau)-a\lambda_E<0\).
\item On \([r_U,r_C]\), write
\(\varphi_\rho=[\rho\DT m_{\lb}-k-a\lb]+\rho\DT[M_{\lb}-\tau]+a[\lb-\lambda_E]\). The first bracket is nonnegative from
\(r_U\) on, and the other two terms are nonnegative up to \(r_C\), with not all zero at once. So \(\varphi_\rho>0\)
there, and its zero \(r_Q\) lies in \((r_N,r_U)\).
\item On \((r_N,r_C]\), \(\lambda_E\le\lb\). The floor-entry purchase posterior falls short of \(\tau\) exactly when
\(\min\{\lambda_Q,\lb\}<\lambda_E\), that is \(\lambda_Q<\lambda_E\). The gap of \eqref{eq:research-floor} falls in
\(\lambda\) and is zero at \(\lambda_Q\), so this is \(\varphi_\rho<0\), that is \(r<r_Q\).
\item Above \(r_C\), \(r>r_U\), so floor research sits at the cap, and \(\tau>M_{\lb}\). Floor-entry research holds
there, uniquely by Lemma A.17.
\end{itemize}
On \([r_N,r_Q]\) the floor root rises from zero to \(\lambda_E(r_Q)\), which is below the cap.

\paragraph{Step 5: nothing else.} Lemma A.16 lists every equilibrium. Above \(r_N\), every equilibrium has positive
research, so it is floor-entry research or recruiting research. With Steps 1 to 4, this gives the five intervals.

\paragraph{Step 6: the six strengths.} At \(r_0\), the gross return of any nonzero order is at most \(\DT(r_0)\), which
is below \(k\). So every such order loses, \(\Psi=0\), research is zero, and \(B(1/2)<c_H\) gives entry \(\rho\). At the
intermediate strengths, \(1/2<\tau<1\), \(\rho\DT<2k\) and \(\lambda_E<\lambda_R<\lb\), so Lemma A.16 leaves exactly zero
research and recruiting research below the cap, with strict low-class abstention. At \(r_1\) and \(r_2\), (U) holds and
\(\tau(r_1)<M_{\lb}<\tau(r_2)\), so Lemma A.17 gives the unique outcome. Class \(H\) outbids every incumbent on the path,
because \(r\le r_2<h\), so its ownership is \(\tfrac12[\rho+(1-\rho)(\beta+(1-\beta)\hat\lambda)]\).

\begin{watchout}
Corollary 3 holds in the example, under the default tie rules, with the paper's research cost and cap. Several of its
steps compare functions at points of the path with exact arithmetic. Proposition~4 holds for every distribution that
meets its hypotheses. Keep the two kinds of claims apart.
\end{watchout}

\subsection{The figure}

\begin{figure}[ht]
\centering
\includegraphics[width=\linewidth]{research}
\caption{Research (panel (a)) and entry (panel (b)) along the example's path, at the prior \(1/2\), under the default
tie rules. Zero research, recruiting research and floor-entry research are drawn as separate curves. Filled markers
belong to a curve and open markers do not. Black dots mark \(r_0\), the three intermediate strengths, \(r_1\) and
\(r_2\). The inset of panel (b) enlarges entry along the recruiting branch.}
\label{fig:research-path}
\end{figure}

Read Figure~\ref{fig:research-path} from left to right.
\begin{itemize}
\item \emph{Zero research} lies flat at zero in panel (a), and its entry is \(\rho\) in panel (b). It runs from \(r_0\)
to \(r_N\). The filled marker at \(r_N\) belongs to it.
\item \emph{Recruiting research} starts at \(r_R\) at a positive level, \(\lambda_E(r_R)\). It does not start at zero,
because a purchase recruits only at research of at least \(\lambda_E\). It rises to the cap at \(r_{\lb}\) and stays
there up to \(r_C\), where the filled marker belongs to it. Its entry rises over the same stretch, from a level well
above \(\rho\), and stays flat at the cap. The inset shows this rise.
\item \emph{Floor-entry research} starts at zero at \(r_N\) and rises toward \(\lambda_E(r_Q)\), below the cap. The open
marker at \(r_Q\) shows that \(r_Q\) itself is excluded. Its entry stays at \(\rho\). Above \(r_C\), floor-entry
research sits at the cap in panel (a) and at \(\rho\) in panel (b); in the figure this last stretch is short.
\item \emph{The pattern.} Two equilibria coexist on \([r_R,r_N]\) and on \((r_N,r_Q)\). On \([r_Q,r_C]\) only recruiting
research remains. Entry first rises and later falls along the path, as in Proposition~2.
\end{itemize}

\paragraph{Comparison with Chapter~\ref{ch:coexistence}.} The recruiting branch starts later than the informative
branch of Chapter~\ref{ch:coexistence}, which starts at \(r_A\), the root of \(J=k\). With research, a purchase must pay
for research as well as for trade: \(J=k+2a\lambda_E\). At \(r_A\), \(\varphi_a=-2a\lambda_E<0\), and \(\varphi_a\) rises,
so \(r_A<r_R\). The two branches also adjust in different ways. With a fixed presence probability, the purchase rate
adjusts through a mixture of the high class. With research, the intensity adjusts, and the high class buys with
probability one after a success.

% =====================================================================
\section{The research cap, the cost and the tie rules}\label{sec:research-ties}

\subsection{Proposition A.1: the research cap and cost}

The paper asks which caps and costs give the pattern of Corollary~3 at its six strengths. It keeps the example's other
primitives and the example's \(\beta\).

\begin{paperresult}{Proposition A.1 (The research cap and cost), in symbols}
Keep the example's primitives other than \(a\) and \(\lb\). Under an indifferent investor who may mix and an
indifferent high-cost challenger who prepares, the six-strength equilibrium pattern of Corollary~3, with
\(\lambda_E(r)<\lambda_R(r)<\lb\) at each intermediate strength, holds with strict margins exactly when
\[
\lambda_E(r_1)<\lb<\lambda_E(r_2),\qquad a_{\rm lo}(\lb)<a<a_{\rm up},
\]
where \(a_{\rm lo}(\lb)=[\DT m_{\lb}-k]/(2\lb)\) at the highest intermediate strength, and
\(a_{\rm up}=[\DT(1-\tau)-k]/(2\lambda_E)\) at the lowest intermediate strength. Entry at \(r_1\) is
\(\rho+(1-\rho)[\beta+(1/2-\beta)\lb]\). With \(\beta\) varying, \(1+u=\eu\), \(\lb=u\beta/(1+u\beta)\) and the paper's
\(a\), the pattern holds for \(\beta\in(\beta_{\rm lo},1/3]\), with a lower bound \(\beta_{\rm lo}\) that the paper
computes.
\end{paperresult}

Each of the four bounds keeps one feature of the pattern.
\begin{itemize}
\item \(\lb>\lambda_E(r_1)\): research at the cap lets a purchase recruit against the strong incumbent,
\(M_{\lb}>\tau(r_1)\). Otherwise entry above \(\rho\) at \(r_1\) would need a tie.
\item \(\lb<\lambda_E(r_2)\): even research at the cap leaves a purchase short of the threshold against the very strong
incumbent, \(M_{\lb}<\tau(r_2)\). Otherwise a purchase would recruit at \(r_2\).
\item \(a<a_{\rm up}\): research is cheap enough that \(\varphi_a>0\) at the lowest intermediate strength. Since
\(\varphi_a\) rises there, it is positive at the other two. Otherwise recruiting research below the cap would be absent
there, or would need a challenger tie.
\item \(a>a_{\rm lo}(\lb)\): research is costly enough that it stays below the cap at the highest intermediate strength.
Since \(\lambda_R\) rises with strength, it stays below the cap at the other two.
\end{itemize}
The cost window is not empty anywhere on the strip of caps. With \(\lb\) tied to \(\beta\) by \(\lb=u\beta/(1+u\beta)\),
\(m_{\lb}=1/(1+\eu)\) does not depend on \(\beta\). So the two cap conditions hold for every \(\beta\), and only the cost
window moves with \(\beta\): a larger \(\beta\) raises \(\lb\) and lowers \(a_{\rm lo}\). Chapter~\ref{app:numbers}
gives the strip, the window and \(\beta_{\rm lo}\).

\subsection{Remark A.1: tie rules with research}

Proposition~4 and Corollary~3 use the default tie rules. The paper's remark says what the other rules add.

\begin{paperresult}{Remark A.1 (Tie rules with research)}
Maintain the general setting. When an indifferent investor may mix and an indifferent high-cost challenger may prepare
with any probability, the outcomes above are supplemented by the following threshold choices. For \(1/2<\tau<1\) and
\(0<\lambda_E(\tau)<\lb\), the threshold pair has research \(\lambda_E\), full high-class purchasing and the low-class
sale rule of Lemma A.15. Its purchase entry is
\[
e^*=\frac{2a\lambda_E+k-(\rho J-k)^+}{J},
\]
and it exists exactly when \(\rho\le e^*\le1\). Entry at a sale and no order is \(\rho\), and the high-cost preparation
probability at a purchase is \((e^*-\rho)/(1-\rho)\). At \(e^*=1\) it coincides with recruiting research; at
\(e^*=\rho\) it is floor-entry research with a purchase posterior equal to \(\tau\).

At \(\lambda_E(\tau)=\lb\), the complete threshold family has purchase entry
\(e(1)\in\bigl[\max\{\rho,\,(2a\lb+k-(\rho J-k)^+)/J\},\,1\bigr]\), provided the lower endpoint is at most one. Each
choice has its own competitive purchase price, research at the cap, and total entry
\(\rho+(e(1)-\rho)[\beta+(1/2-\beta)\lb]\).

When an indifferent investor abstains, a positive-research candidate at intensity \(\lambda\), including a threshold
pair, survives exactly when \(\rho\DT\le2k\) or \(\rho\DT m_\lambda>k\). At zero research the abstention rule selects
abstention after every hypothetical success. Research itself is a unique best response at fixed prices and
preparation, because its objective is strictly concave.
\end{paperresult}

\paragraph{The threshold pair.} Positive research has a strict full purchase (Lemma A.14). Partial preparation can
occur only at a purchase posterior equal to \(\tau\), so \(\lambda=\lambda_E(\tau)\), and \(1-M_{\lambda_E}=1-\tau\).
The purchase then earns \(e(1)\DT(1-\tau)-k=e(1)J-k\).

\begin{derivation}[the purchase entry of the threshold pair]
\begin{align*}
2a\lambda_E&=e(1)\,J-k+(\rho J-k)^+ \why{research condition, Lemma A.12, interior}\\
e(1)\,J&=2a\lambda_E+k-(\rho J-k)^+ \why{add \(k-(\rho J-k)^+\) to both sides}\\
e(1)&=\frac{2a\lambda_E+k-(\rho J-k)^+}{J}=e^* \why{divide by \(J>0\)}
\end{align*}
\end{derivation}

Entry \(e^*=\rho+(1-\rho)\zeta\) needs \(\zeta=(e^*-\rho)/(1-\rho)\in[0,1]\), that is \(\rho\le e^*\le1\). At \(e^*=1\) the
equation is the recruiting equation \eqref{eq:research-cond} at \(\lambda_E\). At \(e^*=\rho\), with \(\rho J>k\), it reads
\(2a\lambda_E=2(\rho J-k)\), the floor equation \eqref{eq:research-floor} at \(\lambda_E\). At the cap, research optimality
asks only that the left side be at most the right side, and this gives the interval of purchase entry.

\paragraph{An investor who abstains.} Suppose \(\rho\DT/2>k\ge\rho\DT m_\lambda\). Then Lemma A.15 needs class \(L\) to
mix with a sale that breaks even. An investor who abstains at a tie cannot mix: with no sales, the sale would pay
strictly; with any sales, it would break even or lose, and the rule then picks abstention. So no low-class rate
qualifies, and the candidate fails. This gives the survival condition.

\paragraph{The example.} The paper's version for the example adds three facts.
\begin{itemize}
\item At the three intermediate strengths, the alternative challenger rule adds the threshold pair, with entry between
\(\rho\) and the recruiting level. At \(r_C\), the cap family permits every \(e(1)\in[\rho,1]\), and total entry fills
the interval from \(\rho\) to the cap level of recruiting research.
\item When an indifferent investor abstains, recruiting research is absent on \((r_N,r_S]\), where the low class would
have to mix, and floor-entry research is the unique equilibrium there.
\item At every strength above \(r_U\) other than \(r_C\), the equilibrium is unique under all four tie rules, by
Lemma A.17.
\end{itemize}

\begin{watchout}
State the tie rule of each result. Lemmas A.12, A.13 and A.17 hold under all four tie rules. Lemma A.14 holds under
either investor rule and either challenger rule. Lemma A.15, Lemma A.16, Proposition~4 and Corollary~3 use the
default tie rules. Remark A.1 lists what the other rules add. Research itself never has a tie.
\end{watchout}

\subsection{Access to prices with research}

The appendix on research also compares two economies under condition (U) with \(\tau<M_{\lb}\): in one the challenger
sees the price, and in the other it does not. Chapter~\ref{ch:welfare} studies access to prices in general. Under all
four tie rules, both economies have full orders and research at the cap, uniquely. With access, only a purchase
recruits. With the price hidden, the challenger keeps its prior \(1/2<\tau\), and entry is \(\rho\) after every order.
So the order law, the research cost and the trading cost are the same in both.

Access adds high-cost preparation after a purchase. A purchase has probability \(\bar b+\lb/2\), with
\(\bar b=(1-\lb)\beta\). The added preparation probability is \(1-\rho\). The posterior after a purchase is \(M_{\lb}\),
so each added entry raises expected proceeds by \(w(M_{\lb})=t_L-t_0+\DT M_{\lb}\). Multiply:
\[
\text{proceeds gain}=(1-\rho)\,(\bar b+\lb/2)\,\bigl[t_L-t_0+\DT M_{\lb}\bigr]>0 .
\]
For surplus, an added prepared challenger adds \(B(M_{\lb})+pF(p)\) to the expected value of the allocation, and costs
\(c_H\). Since \(B(M_{\lb})-c_H=\gamma(M_{\lb}-\tau)\), the paper obtains
\[
\text{surplus gain}=(1-\rho)\,(\bar b+\lb/2)\,\bigl[\gamma(M_{\lb}-\tau)+pF(p)\bigr]>0 .
\]
Both are positive because \(M_{\lb}>\tau\), \(\gamma>0\), \(\DT>0\) and \(t_L\ge t_0\).

\begin{chaptersummary}
Research is the presence probability chosen at a cost \((a/2)\lambda^2\), up to a cap \(\lb<1\); a success fills the
slot, and a failure passes it to a liquidity trader. At the prior \(1/2\), the top posterior is
\(M_\lambda=1-\beta(1-\lambda)/[2\beta+(1-2\beta)\lambda]\); it rises with research, and a full purchase recruits
exactly when research is at least \(\lambda_E(\tau)\). In recruiting research, the value of a success gives the
condition \(2a\lambda=\DT m_\lambda-k+(\rho\DT m_\lambda-k)^+\), whose left side rises and right side falls, so research
is \(\min\{\lambda_R(\DT),\lb\}\) when \(\DT>2k\), with no tie. Under the default tie rules, Proposition~4 says that
research rises with the spread up to the cap, entry \(\rho+(1-\rho)[\beta+(1/2-\beta)\lambda]\) and class-\(H\) entry
rise, and class-\(L\) entry falls. High-class ownership is half the entry in class \(H\) when class \(H\) outbids the
incumbent. When the exact recruitment margin rises, the equilibrium holds on one interval of strengths. In the
example, zero research, recruiting research and floor-entry research divide the path at \(r_R\), \(r_N\), \(r_Q\) and
\(r_C\), and entry rises and later falls.
\end{chaptersummary}

% =====================================================================
\needspace{10\baselineskip}
\answersection

\answer{\ref{q:research-lamE}}{\emph{A second derivation.} Set \(m_\lambda=1-\tau\) with \eqref{eq:research-bounds}:
\begin{align*}
\beta(1-\lambda)&=(1-\tau)\,[2\beta+(1-2\beta)\lambda] \why{multiply by the positive denominator}\\
\beta-\beta\lambda&=2\beta(1-\tau)+(1-2\beta)(1-\tau)\lambda \why{expand both sides}\\
\beta-2\beta(1-\tau)&=\lambda\,[\beta+(1-2\beta)(1-\tau)] \why{move the \(\lambda\) terms to the right}\\
\beta(2\tau-1)&=\lambda\,[\beta+(1-2\beta)(1-\tau)] \why{\(\beta-2\beta+2\beta\tau=\beta(2\tau-1)\)}\\
\lambda&=\frac{\beta(2\tau-1)}{\beta+(1-2\beta)(1-\tau)} \why{divide by the positive bracket}
\end{align*}
This is \eqref{eq:research-lamE}.

\emph{The effect of \(\beta\).} Write \(c=2\tau-1>0\) and \(d=1-\tau>0\). The denominator is
\(\beta+(1-2\beta)d=d+\beta(1-2d)=d+\beta c\), since \(1-2d=2\tau-1=c\). So \(\lambda_E=\beta c/(d+\beta c)\), and
\begin{align*}
\frac{\partial\lambda_E}{\partial\beta}&=\frac{c\,(d+\beta c)-\beta c\cdot c}{(d+\beta c)^2} \why{quotient rule}\\
&=\frac{cd}{(d+\beta c)^2}=\frac{(2\tau-1)(1-\tau)}{[\beta+(1-2\beta)(1-\tau)]^2}>0 \why{the \(\beta c^2\) terms cancel}
\end{align*}
In words: a larger \(\beta\) means more liquidity orders in each direction, so the investor's orders have more
camouflage. The same research then moves the posterior after a purchase less, and more research is needed before a
purchase reaches \(\tau\).}

\answer{\ref{q:research-zero}}{At \(\hat\lambda=0\), the order probabilities are \(f_j(x)=(1-0)\eta(x)+0\cdot\sigma_j(x)
=\eta(x)\) in both classes, whatever \(\sigma_j\) is. Bayes' rule at the prior \(1/2\) gives
\[
\mu(x)=\frac{\tfrac12\eta(x)}{\tfrac12\eta(x)+\tfrac12\eta(x)}=\frac12
\]
after every order. A belief at a price averages these posteriors, so it is \(1/2\) too. Since \(\tau>1/2\), a high-cost
challenger strictly prefers not to prepare, under either challenger rule, and entry is \(\rho\) after every order.

After a hypothetical success, the correct orders earn
\[
U_H(1)=\rho\DT\bigl(1-\tfrac12\bigr)-k=\tfrac12\rho\DT-k,\qquad U_L(-1)=\rho\DT\cdot\tfrac12-k=\tfrac12\rho\DT-k,
\]
and abstention earns zero. Wrong-signed orders earn \(-\rho\DT/2-k<0\). So
\(\Psi=\tfrac12(\tfrac12\rho\DT-k)^++\tfrac12(\tfrac12\rho\DT-k)^+=(\tfrac12\rho\DT-k)^+\). Lemma A.12 gives the best
research \(\min\{\lb,\Psi/a\}\). It is zero exactly when \(\Psi=0\), that is \(\rho\DT\le2k\). If \(\rho\DT>2k\), the
best research is positive, and zero research fails (E7).

The orders after a hypothetical success must be optimal because they set the value of research. A deviation in research
makes a success possible. Without the requirement, one could support zero research with poor orders after a success,
which would make research look worthless. The paper therefore asks (E5) to hold at \(\hat\lambda=0\).}

\answer{\ref{q:research-floor}}{In floor-entry research, no order recruits, so entry is \(\rho\) after every order.
Full orders give \(\mu(1)=M_\lambda\) and \(\mu(-1)=m_\lambda\). By \eqref{eq:trading-returns},
\begin{align*}
U_H(1)&=\rho\DT(1-M_\lambda)-k=\rho\DT m_\lambda-k \why{\(e(1)=\rho\); \(1-M_\lambda=m_\lambda\)}\\
U_L(-1)&=\rho\DT m_\lambda-k \why{\(e(-1)=\rho\), \(\mu(-1)=m_\lambda\)}
\end{align*}
Both are the best orders, so \(\Psi=\tfrac12(\rho\DT m_\lambda-k)+\tfrac12(\rho\DT m_\lambda-k)=\rho\DT m_\lambda-k\).
Below the cap, Lemma A.12 gives \(a\lambda=\Psi\), which is \eqref{eq:research-floor}.

The factors differ because of what a success earns in each class. In floor-entry research both classes earn the same
amount after a success, so the average over the two classes is that amount itself. In recruiting research the two
classes earn different amounts, \(\DT m_\lambda-k\) and \((\rho\DT m_\lambda-k)^+\). Their average carries the factor
\(\tfrac12\), and multiplying the condition by \(2\) moves it to the left side as \(2a\lambda\).}

\answer{\ref{q:research-entry}}{Entry in class \(L\) is \(e_L=\rho+(1-\rho)\beta(1-\hat\lambda)\). Class \(L\) never buys,
so it reaches a purchase, the only order that recruits, only through a liquidity trader, with probability
\(b=(1-\hat\lambda)\beta\). Research takes the slot away from liquidity traders, so this probability falls, and with it
the chance that the market mistakes class \(L\) for class \(H\). So \(e_L\) falls at the rate \((1-\rho)\beta\).

Total entry has the slope \((1-\rho)(1/2-\beta)\). At \(\beta=1/2\), this slope is zero: each unit of research would add
an informed purchase half the time and remove a liquidity purchase half the time, and the two changes would cancel. The
model excludes \(\beta=1/2\), because then a liquidity trader never abstains, \(b_0=(1-\lambda)(1-2\beta)=0\), and the
zero order would have no liquidity probability. The research results take \(0<\beta<1/2\).}

\answer{\ref{q:research-kink}}{Section~\ref{sec:research-sale} gives the cap spread in the case
\(\lb\le(1-\rho)k/(2a\rho)\): \(\DT=(2a\lb+k)/m_{\lb}\). At this spread,
\begin{align*}
\rho\DT m_{\lb}&=\rho(2a\lb+k) \why{substitute the cap spread}\\
&\le\rho\Bigl(\frac{(1-\rho)k}{\rho}+k\Bigr)=k \why{\(2a\lb\le(1-\rho)k/\rho\)}
\end{align*}
So a sure sale does not pay strictly when research reaches the cap. Below the cap spread, the spread is smaller, so
\(\rho\DT m_\lambda\) at the research level is at most \(k\) as well (research and the spread move together, and the
sale boundary meets the first piece only at the intensity \((1-\rho)k/(2a\rho)\ge\lb\)). So the sale term is zero at
every uncapped root. Beyond the cap spread, research sits at the cap and cannot rise.

In words: by the time a sure sale would add a second return to research, research has no room left to grow. The sale's
profit then changes what the investor earns, and leaves its research choice unchanged. The example has this case: there
the ratio \((1-\rho)k/(2a\rho)\) exceeds one.}
